% oss-plantilla: oss-2026.11
% Preámbulo corporativo fijo. No editar en sd-NN.tex: exportar_latex.py verificar lo compara byte a byte.
% Para cambiar el formato, editar esta plantilla u oss.sty y subir la versión de la primera línea.
\PassOptionsToPackage{unicode}{hyperref}
\PassOptionsToPackage{hyphens}{url}
\documentclass[11pt]{article}
\usepackage{amsmath,amssymb}
\usepackage{iftex}
\ifPDFTeX
  % Prism y otros editores en la nube: pdfLaTeX con Helvetica como familia por defecto.
  \usepackage[T1]{fontenc}
  \usepackage[utf8]{inputenc}
  \usepackage{textcomp}
  \usepackage[scaled=0.92]{helvet}
  \renewcommand{\familydefault}{\sfdefault}
  \DeclareUnicodeCharacter{2264}{\ensuremath{\leq}}
  \DeclareUnicodeCharacter{2265}{\ensuremath{\geq}}
  \DeclareUnicodeCharacter{2260}{\ensuremath{\neq}}
  \DeclareUnicodeCharacter{2248}{\ensuremath{\approx}}
  \DeclareUnicodeCharacter{2190}{\ensuremath{\leftarrow}}
  \DeclareUnicodeCharacter{2192}{\ensuremath{\rightarrow}}
  \DeclareUnicodeCharacter{25CF}{\ensuremath{\bullet}}
  \DeclareUnicodeCharacter{2713}{\ensuremath{\checkmark}}
\else
  % XeLaTeX/LuaLaTeX (compilación oficial local): DejaVu Sans.
  \usepackage{fontspec}
  \defaultfontfeatures{Ligatures=TeX,Scale=MatchLowercase}
  \setmainfont{DejaVu Sans}
  \setsansfont{DejaVu Sans}
  \setmonofont{DejaVu Sans Mono}
\fi
\usepackage[letterpaper,margin=20mm]{geometry}
\usepackage{microtype}
\usepackage{parskip}
\usepackage{xcolor}
\usepackage{longtable,booktabs,array}
\usepackage{calc}
\usepackage{etoolbox}
\makeatletter
\patchcmd\longtable{\par}{\if@noskipsec\mbox{}\fi\par}{}{}
\makeatother
\usepackage{footnotehyper}
\makesavenoteenv{longtable}
\usepackage{graphicx}
\makeatletter
\newsavebox\pandoc@box
\newcommand*\pandocbounded[1]{%
  \sbox\pandoc@box{#1}%
  \Gscale@div\@tempa{\textheight}{\dimexpr\ht\pandoc@box+\dp\pandoc@box\relax}%
  \Gscale@div\@tempb{\linewidth}{\wd\pandoc@box}%
  \ifdim\@tempb\p@<\@tempa\p@\let\@tempa\@tempb\fi
  \ifdim\@tempa\p@<\p@\scalebox{\@tempa}{\usebox\pandoc@box}%
  \else\usebox{\pandoc@box}%
  \fi}
\def\maxwidth{\ifdim\Gin@nat@width>\linewidth\linewidth\else\Gin@nat@width\fi}
\def\maxheight{\ifdim\Gin@nat@height>\textheight\textheight\else\Gin@nat@height\fi}
\makeatother
\setkeys{Gin}{width=\maxwidth,height=\maxheight,keepaspectratio}
\makeatletter
\def\fps@figure{htbp}
\makeatother
\usepackage[normalem]{ulem}
\providecommand{\st}[1]{\sout{#1}}
\usepackage{fancyvrb}
\setlength{\emergencystretch}{3em}
\providecommand{\tightlist}{%
  \setlength{\itemsep}{0pt}\setlength{\parskip}{0pt}}
\providecommand{\passthrough}[1]{#1}
\setcounter{secnumdepth}{-\maxdimen}
\ifPDFTeX
  \usepackage[spanish]{babel}
\else
  \usepackage[bidi=default]{babel}
  \babelprovide[main,import]{spanish}
\fi
\let\LanguageShortHands\languageshorthands
\def\languageshorthands#1{}
\usepackage{bookmark}
\usepackage{xurl}
\urlstyle{same}
\hypersetup{pdflang={es-CL},pdfcreator={Only Simple Solutions - exportar_latex.py}}
\usepackage{oss}
\begin{document}

\hypersetup{pageanchor=false}
\ossCover{Formulario T-9}{Subdocumento 6}
\hypersetup{pageanchor=true}
\tableofcontents
\clearpage

\renewcommand{\tablename}{Tabla}
\renewcommand{\thetable}{T-9.\arabic{table}}

\section{Formulario T-9 · Metodología para la administración y gestión del proyecto}\label{formulario-t-9}

Conforme al Formulario T-7 y al Formulario T-21 de las Bases Administrativas, este formulario adjunta el detalle de la información solicitada en el Subdocumento 6, letra a, Metodología de gestión de proyectos. La sección 6.1 del Subdocumento 6 contiene el resumen y el análisis; este formulario reúne los listados, las definiciones operativas, los procedimientos y las plantillas con que Only Simple Solutions dirigirá el proyecto de Multitiendas Ancoa S.A.

El formulario sigue el orden de la sección 6.1. La Tabla~\mbox{\ref{tab:t9-indice}} indica qué contiene cada apartado y con qué sección del Subdocumento 6 se corresponde.

\begin{longtable}[]{@{}
  >{\raggedright\arraybackslash}p{(\columnwidth - 6\tabcolsep) * \real{0.3000}}
  >{\raggedright\arraybackslash}p{(\columnwidth - 6\tabcolsep) * \real{0.5000}}
  >{\raggedright\arraybackslash}p{(\columnwidth - 6\tabcolsep) * \real{0.2000}}@{}}
\caption{Correspondencia entre el formulario y la sección 6.1.}\label{tab:t9-indice}\\
\toprule\noalign{}
\begin{minipage}[b]{\linewidth}\raggedright
\textbf{Apartado}
\end{minipage} & \begin{minipage}[b]{\linewidth}\raggedright
\textbf{Contenido}
\end{minipage} & \begin{minipage}[b]{\linewidth}\raggedright
\textbf{Sección del Subdocumento 6}
\end{minipage} \\
\midrule\noalign{}
\endfirsthead
\toprule\noalign{}
\begin{minipage}[b]{\linewidth}\raggedright
\textbf{Apartado}
\end{minipage} & \begin{minipage}[b]{\linewidth}\raggedright
\textbf{Contenido}
\end{minipage} & \begin{minipage}[b]{\linewidth}\raggedright
\textbf{Sección del Subdocumento 6}
\end{minipage} \\
\midrule\noalign{}
\endhead
\bottomrule\noalign{}
\multicolumn{3}{@{}l}{\rule{0pt}{3.2ex}\textbf{Fuente:} Elaboración propia.} \\
\endlastfoot
\textbf{T-9.1 Adaptación del PMBOK} & Procesos aplicados por área de conocimiento, enfoque por tipo de trabajo y flujos Kanban. & 6.1.1 \\ \addlinespace
\textbf{T-9.2 Gobierno, roles y decisiones} & Roles y límites de autoridad, instancias de gobierno, escalamiento y matriz de decisiones. & 6.1.2 y 6.1.7 \\ \addlinespace
\textbf{T-9.3 Integración y control de cambios} & Siete procesos de integración, procedimiento de cambios, valor ganado y condiciones de cierre. & 6.1.3 \\ \addlinespace
\textbf{T-9.4 Gestión de interesados} & Registro de los 26 interesados, matriz de involucramiento, dueños de producto, tensiones, ADKAR e indicadores. & 6.1.4 \\ \addlinespace
\textbf{T-9.5 Gestión de comunicaciones} & Sistema de información, matriz completa, contenidos mínimos y monitoreo. & 6.1.5 \\ \addlinespace
\textbf{T-9.6 Gestión de adquisiciones} & Registro de dependencias, catálogo por tipo, estados, responsabilidades y cierre de terceros. & 6.1.6 \\ \addlinespace
\textbf{T-9.7 Gobierno por etapas} & Énfasis del gobierno por período, desde el mes 1 hasta el mes 56. & 6.1.7.1 \\ \addlinespace
\textbf{T-9.8 Plantillas} & Solicitud de cambio, acta de comité y lección aprendida. & 6.1.2 y 6.1.3 \\
\end{longtable}

Todos los registros que se describen se mantendrán en el sistema de información para la dirección del proyecto (PMIS), formado por Jira y Microsoft SharePoint. Este formulario no contiene precios, tarifas ni montos que permitan inferir la Oferta Económica.

\subsection{T-9.1 Adaptación del PMBOK y enfoque de trabajo}\label{t9-1}

Only Simple Solutions aplicará la Guía del PMBOK, sexta edición (PMI, 2017), conservando los procesos necesarios para controlar los compromisos de esta solución. La Tabla~\mbox{\ref{tab:t9-procesos}} relaciona cada área de conocimiento con los procesos que se aplican, su responsable y el producto que deja evidencia de su ejecución.

\begin{longtable}[]{@{}
  >{\raggedright\arraybackslash}p{(\columnwidth - 8\tabcolsep) * \real{0.1900}}
  >{\raggedright\arraybackslash}p{(\columnwidth - 8\tabcolsep) * \real{0.3200}}
  >{\raggedright\arraybackslash}p{(\columnwidth - 8\tabcolsep) * \real{0.1900}}
  >{\raggedright\arraybackslash}p{(\columnwidth - 8\tabcolsep) * \real{0.3000}}@{}}
\caption{Procesos aplicados por área de conocimiento.}\label{tab:t9-procesos}\\
\toprule\noalign{}
\begin{minipage}[b]{\linewidth}\raggedright
\textbf{Área}
\end{minipage} & \begin{minipage}[b]{\linewidth}\raggedright
\textbf{Procesos aplicados}
\end{minipage} & \begin{minipage}[b]{\linewidth}\raggedright
\textbf{Responsable}
\end{minipage} & \begin{minipage}[b]{\linewidth}\raggedright
\textbf{Producto y evidencia}
\end{minipage} \\
\midrule\noalign{}
\endfirsthead
\toprule\noalign{}
\begin{minipage}[b]{\linewidth}\raggedright
\textbf{Área}
\end{minipage} & \begin{minipage}[b]{\linewidth}\raggedright
\textbf{Procesos aplicados}
\end{minipage} & \begin{minipage}[b]{\linewidth}\raggedright
\textbf{Responsable}
\end{minipage} & \begin{minipage}[b]{\linewidth}\raggedright
\textbf{Producto y evidencia}
\end{minipage} \\
\midrule\noalign{}
\endhead
\bottomrule\noalign{}
\multicolumn{4}{@{}l}{\rule{0pt}{3.2ex}\textbf{Fuente:} Elaboración propia a partir de PMI (2017).} \\
\endlastfoot
\textbf{Integración} & Los siete procesos: acta de constitución, plan para la dirección, dirigir y gestionar el trabajo, gestionar el conocimiento, monitorear y controlar, control integrado de cambios y cierre del proyecto o fase. & Jefe de Proyecto & Acta de constitución, plan para la dirección aprobado en H1, registro de cambios, registro de lecciones aprendidas y actas de cierre. \\ \addlinespace
\textbf{Alcance} & Planificar, recopilar requisitos, definir el alcance, crear la EDT, validar y controlar el alcance. & Jefe de Proyecto con el Líder Funcional & Enunciado del alcance, EDT y su diccionario, línea base de alcance y matriz de trazabilidad (Formulario T-12). \\ \addlinespace
\textbf{Cronograma} & Planificar, definir y secuenciar actividades, estimar duraciones, desarrollar y controlar el cronograma. & Jefe de Proyecto & Línea base de cronograma, carta Gantt de 56 meses (Formulario T-14) y pronóstico actualizado. \\ \addlinespace
\textbf{Costos} & Planificar, estimar, determinar el presupuesto y controlar los costos. & Jefe de Proyecto & Línea base de costos y medición de valor ganado, comunicadas por los canales contractuales. \\ \addlinespace
\textbf{Calidad} & Planificar, gestionar y controlar la calidad. & Líder de Calidad y Pruebas & Plan de calidad (Subdocumento 9), puertas de calidad e informes de pruebas (Formulario T-13). \\ \addlinespace
\textbf{Recursos} & Planificar, estimar, adquirir, desarrollar, dirigir y controlar los recursos. & Jefe de Proyecto & Asignaciones y suplencias (Formulario T-8) y nivelación de recursos (Formulario T-15). \\ \addlinespace
\textbf{Comunicaciones} & Planificar, gestionar y monitorear las comunicaciones. & Jefe de Proyecto & Plan de comunicaciones aprobado en H1, matriz, actas, informe mensual y bitácora de monitoreo. \\ \addlinespace
\textbf{Riesgos} & Planificar, identificar, analizar de forma cualitativa y cuantitativa, planificar e implementar respuestas y monitorear, conforme a ISO 31000. & Jefe de Proyecto con los líderes de frente & Registro vivo de riesgos revisado en cada Comité de Proyecto (Subdocumento 8 y Formulario T-16). \\ \addlinespace
\textbf{Adquisiciones} & Planificar, efectuar y controlar las adquisiciones y las dependencias externas. & Jefe de Proyecto con el Arquitecto de Solución & Registro de dependencias en Jira y actas de recepción técnica. \\ \addlinespace
\textbf{Interesados} & Identificar, planificar el involucramiento, gestionar la participación y monitorear el involucramiento. & Jefe de Proyecto con el Líder de Implantación y Gestión del Cambio & Registro de interesados, estrategia de involucramiento aprobada en H1 y matriz de evaluación. \\
\end{longtable}

Los planes subsidiarios forman parte de un solo plan para la dirección del proyecto. Cuando cambie una dependencia, el Jefe de Proyecto revisará en conjunto los planes afectados, en lugar de actualizarlos por separado.

\subsubsection{Enfoque por tipo de trabajo}\label{t9-1-enfoque}

La metodología es híbrida. La Tabla~\mbox{\ref{tab:t9-enfoque}} asigna a cada tipo de trabajo su enfoque, su cadencia y el control que limita lo que ese enfoque puede decidir.

\begin{longtable}[]{@{}
  >{\raggedright\arraybackslash}p{(\columnwidth - 8\tabcolsep) * \real{0.2700}}
  >{\raggedright\arraybackslash}p{(\columnwidth - 8\tabcolsep) * \real{0.1500}}
  >{\raggedright\arraybackslash}p{(\columnwidth - 8\tabcolsep) * \real{0.2600}}
  >{\raggedright\arraybackslash}p{(\columnwidth - 8\tabcolsep) * \real{0.3200}}@{}}
\caption{Enfoque de gestión por tipo de trabajo.}\label{tab:t9-enfoque}\\
\toprule\noalign{}
\begin{minipage}[b]{\linewidth}\raggedright
\textbf{Tipo de trabajo}
\end{minipage} & \begin{minipage}[b]{\linewidth}\raggedright
\textbf{Enfoque}
\end{minipage} & \begin{minipage}[b]{\linewidth}\raggedright
\textbf{Cadencia}
\end{minipage} & \begin{minipage}[b]{\linewidth}\raggedright
\textbf{Control y límite}
\end{minipage} \\
\midrule\noalign{}
\endfirsthead
\toprule\noalign{}
\begin{minipage}[b]{\linewidth}\raggedright
\textbf{Tipo de trabajo}
\end{minipage} & \begin{minipage}[b]{\linewidth}\raggedright
\textbf{Enfoque}
\end{minipage} & \begin{minipage}[b]{\linewidth}\raggedright
\textbf{Cadencia}
\end{minipage} & \begin{minipage}[b]{\linewidth}\raggedright
\textbf{Control y límite}
\end{minipage} \\
\midrule\noalign{}
\endhead
\bottomrule\noalign{}
\multicolumn{4}{@{}l}{\rule{0pt}{3.2ex}\textbf{Fuente:} Elaboración propia.} \\
\endlastfoot
\textbf{Líneas base, hitos, aceptación, cambios y relación contractual} & Predictivo (PMBOK) & Comité de Proyecto quincenal y Comité Ejecutivo mensual & Toda modificación de una línea base requiere solicitud formal y aprobación previa del Comité Ejecutivo. \\ \addlinespace
\textbf{Construcción de servicios e integraciones} & Scrum & Sprints de dos semanas; intervalo de coordinación de cuatro sprints (ocho semanas) & Definición de terminado. La revisión del sprint no sustituye la aceptación contractual ni autoriza producción. \\ \addlinespace
\textbf{Solicitudes de ambientes, tareas de plataforma e incidencias} & Kanban & Flujo continuo; revisión semanal de bloqueos y antigüedad & Límites de trabajo en curso. No incorpora nuevos compromisos de alcance. \\ \addlinespace
\textbf{Adquisiciones, instalaciones, migraciones y pasos a producción} & Predictivo con condiciones de entrada & Programación integrada & El trabajo dependiente no se inicia sin verificar sus condiciones de entrada. \\
\end{longtable}

La cadencia de coordinación de ocho semanas se relacionará con los meses e hitos contractuales en el calendario integrado. No crea compromisos de alcance ni modifica fechas contractuales.

\subsubsection{Flujos Kanban}\label{t9-1-kanban}

Los responsables de cada flujo definirán qué constituye una solicitud, cuándo empieza y termina, sus estados, sus políticas de prioridad y el control del trabajo en curso (Kanban Guides, 2025). La Tabla~\mbox{\ref{tab:t9-kanban}} define los tres flujos del proyecto.

\begin{longtable}[]{@{}
  >{\raggedright\arraybackslash}p{(\columnwidth - 8\tabcolsep) * \real{0.1800}}
  >{\raggedright\arraybackslash}p{(\columnwidth - 8\tabcolsep) * \real{0.3400}}
  >{\raggedright\arraybackslash}p{(\columnwidth - 8\tabcolsep) * \real{0.2000}}
  >{\raggedright\arraybackslash}p{(\columnwidth - 8\tabcolsep) * \real{0.2800}}@{}}
\caption{Flujos Kanban del proyecto.}\label{tab:t9-kanban}\\
\toprule\noalign{}
\begin{minipage}[b]{\linewidth}\raggedright
\textbf{Flujo}
\end{minipage} & \begin{minipage}[b]{\linewidth}\raggedright
\textbf{Qué constituye una solicitud}
\end{minipage} & \begin{minipage}[b]{\linewidth}\raggedright
\textbf{Responsable}
\end{minipage} & \begin{minipage}[b]{\linewidth}\raggedright
\textbf{Política de prioridad}
\end{minipage} \\
\midrule\noalign{}
\endfirsthead
\toprule\noalign{}
\begin{minipage}[b]{\linewidth}\raggedright
\textbf{Flujo}
\end{minipage} & \begin{minipage}[b]{\linewidth}\raggedright
\textbf{Qué constituye una solicitud}
\end{minipage} & \begin{minipage}[b]{\linewidth}\raggedright
\textbf{Responsable}
\end{minipage} & \begin{minipage}[b]{\linewidth}\raggedright
\textbf{Política de prioridad}
\end{minipage} \\
\midrule\noalign{}
\endhead
\bottomrule\noalign{}
\multicolumn{4}{@{}l}{\rule{0pt}{3.2ex}\textbf{Fuente:} Elaboración propia.} \\
\endlastfoot
\textbf{Solicitudes de ambientes} & Creación, cambio o acceso a un ambiente de Desarrollo, QA, Preproducción o Recuperación ante Desastres. & Líder de Operación / SRE & Por la fecha en que el trabajo dependiente necesita el ambiente. \\ \addlinespace
\textbf{Tareas de plataforma} & Habilitación de identidad, integración, observabilidad, flujo de entrega o infraestructura como código. & Líder de Desarrollo, con el equipo de plataforma & Por el hito y la dependencia que habilita. \\ \addlinespace
\textbf{Incidencias} & Falla o degradación en un ambiente y, desde la marcha blanca, en el servicio. & Líder de Operación / SRE & Por severidad. Los tiempos contractuales de atención prevalecen sobre cualquier pronóstico. \\
\end{longtable}

Los tres flujos comparten los estados y las reglas siguientes.

\begin{itemize}
\tightlist
\item
  \textbf{Solicitada:} la solicitud está registrada en Jira con solicitante, descripción y fecha requerida.
\item
  \textbf{Preparada:} cumple el criterio de entrada, es decir, identifica el trabajo que condiciona, su responsable y los antecedentes necesarios para comenzar.
\item
  \textbf{En curso:} una persona del flujo la ejecuta. Un bloqueo se marca sobre la tarjeta, con su causa y la fecha en que se detectó.
\item
  \textbf{En verificación:} el solicitante o el responsable de calidad comprueba el resultado.
\item
  \textbf{Terminada:} cumple el criterio de salida, con resultado verificado y evidencia vinculada en Jira.
\end{itemize}

El límite inicial de trabajo en curso de cada estado se fijará en el mes 1, según las personas asignadas al flujo, y se revisará al cierre de cada intervalo con datos de tiempo de ciclo. La expectativa de servicio se expresará como plazo y probabilidad; por ejemplo, el porcentaje de solicitudes terminadas dentro de un número de días hábiles. Se estimará al inicio y se recalibrará con el tiempo de ciclo observado. El responsable del flujo revisará cada semana bloqueos, antigüedad del trabajo y tiempos de resolución. Este pronóstico de gestión no sustituye los tiempos contractuales de atención de incidentes.

\subsection{T-9.2 Gobierno, roles y decisiones}\label{t9-2}

\subsubsection{Roles y límites de autoridad}\label{t9-2-roles}

La Tabla~\mbox{\ref{tab:t9-roles}} reúne los roles de gestión de ambas partes, su responsabilidad principal y el límite de su autoridad. La composición nominal, las dedicaciones y las suplencias se declaran en el Formulario T-8 y en la nivelación del Formulario T-15, sin contar dos veces la capacidad de una persona.

\begin{longtable}[]{@{}
  >{\raggedright\arraybackslash}p{(\columnwidth - 8\tabcolsep) * \real{0.2000}}
  >{\raggedright\arraybackslash}p{(\columnwidth - 8\tabcolsep) * \real{0.1700}}
  >{\raggedright\arraybackslash}p{(\columnwidth - 8\tabcolsep) * \real{0.3500}}
  >{\raggedright\arraybackslash}p{(\columnwidth - 8\tabcolsep) * \real{0.2800}}@{}}
\caption{Roles de gestión y límites de autoridad.}\label{tab:t9-roles}\\
\toprule\noalign{}
\begin{minipage}[b]{\linewidth}\raggedright
\textbf{Rol}
\end{minipage} & \begin{minipage}[b]{\linewidth}\raggedright
\textbf{Parte}
\end{minipage} & \begin{minipage}[b]{\linewidth}\raggedright
\textbf{Responsabilidad principal}
\end{minipage} & \begin{minipage}[b]{\linewidth}\raggedright
\textbf{Límite de autoridad}
\end{minipage} \\
\midrule\noalign{}
\endfirsthead
\toprule\noalign{}
\begin{minipage}[b]{\linewidth}\raggedright
\textbf{Rol}
\end{minipage} & \begin{minipage}[b]{\linewidth}\raggedright
\textbf{Parte}
\end{minipage} & \begin{minipage}[b]{\linewidth}\raggedright
\textbf{Responsabilidad principal}
\end{minipage} & \begin{minipage}[b]{\linewidth}\raggedright
\textbf{Límite de autoridad}
\end{minipage} \\
\midrule\noalign{}
\endhead
\bottomrule\noalign{}
\multicolumn{4}{@{}l}{\rule{0pt}{3.2ex}\textbf{Fuente:} Elaboración propia.} \\
\endlastfoot
\textbf{Patrocinador} & Ancoa (Gerencia General) & Representa las prioridades estratégicas del cliente en el Comité Ejecutivo. & Decide en el Comité Ejecutivo; no reemplaza al administrador del contrato. \\ \addlinespace
\textbf{Administrador del contrato} & Ancoa & Supervisa el cumplimiento contractual, aprueba los estados de pago y gestiona modificaciones y sanciones. & Toda modificación contractual requiere acuerdo escrito de ambas partes. \\ \addlinespace
\textbf{Contraparte técnica} & Ancoa (TI) & Valida entregables, aprueba avances técnicos, coordina pruebas y suscribe las actas de aceptación de hitos. & Es la única instancia que emite la conformidad de un hito y del paso a producción. \\ \addlinespace
\textbf{Contraparte de cumplimiento} & Ancoa (Contraloría y Negocio Financiero) & Evalúa las obligaciones del emisor fiscalizado en las decisiones que afectan la frontera Retail--Emisor. & Su evaluación documentada es previa a toda aprobación de un cruce. \\ \addlinespace
\textbf{Dueño de producto} & Ancoa, uno por producto & Ordena el backlog de su producto dentro del alcance aprobado y valida funcionalmente los incrementos. & No aprueba cambios contractuales, no acepta hitos ni autoriza pasos a producción. \\ \addlinespace
\textbf{Jefe de Proyecto} & OSS, dedicación exclusiva en implementación & Dirige la ejecución, integra los frentes, consolida el desempeño y coordina cambios, comunicaciones y cierres. & Compromete a la empresa en materias de ejecución y eleva lo que excede sus facultades. \\ \addlinespace
\textbf{Arquitecto de Solución} & OSS & Vela por la coherencia entre los componentes de nube y locales y prepara las decisiones de arquitectura. & Las decisiones transversales se aprueban en el Comité de Arquitectura. \\ \addlinespace
\textbf{Encargado de Seguridad de la Información} & OSS & Verifica accesos y segregación y conduce la notificación de incidentes y brechas. & Activa el aviso de incidente crítico sin esperar una sesión de comité. \\ \addlinespace
\textbf{Líder de Datos} & OSS, permanente en implementación & Evalúa autoridad, integridad y tratamiento de la información; conduce migración y conciliación. & Su evaluación es requisito de toda decisión de cruce entre ambos negocios. \\ \addlinespace
\textbf{Líder de Desarrollo} & OSS, dedicación completa en implementación & Coordina a los equipos Scrum y la planificación de cada intervalo. & Actúa dentro de las líneas base y decisiones vigentes. \\ \addlinespace
\textbf{Líder Funcional} & OSS, dedicación completa en implementación & Verifica el alcance y las reglas de negocio con las contrapartes usuarias. & No modifica requisitos sin control integrado de cambios. \\ \addlinespace
\textbf{Líder de Calidad y Pruebas} & OSS & Define puertas de calidad, criterios de recepción y evidencias de verificación. & No incorpora a un entregable un resultado con observaciones bloqueantes. \\ \addlinespace
\textbf{Líder de Integración} & OSS, permanente en implementación & Coordina contratos de interfaz y dependencias con plataformas y proveedores. & Los cambios de contrato de interfaz siguen el registro de decisiones. \\ \addlinespace
\textbf{Líder de Operación / SRE} & OSS, desde el mes 6 & Coordina soporte, monitoreo, continuidad y atención de incidencias que afectan al servicio. & Autoriza técnicamente las liberaciones; la conformidad productiva es de la contraparte técnica. \\ \addlinespace
\textbf{Líder de Implantación y Gestión del Cambio} & OSS, desde el mes 8 & Prepara sitios y tiendas, capacitación y adopción con el modelo ADKAR. & No dispensa capacitación ni certificación exigidas. \\ \addlinespace
\textbf{Scrum Master} & OSS, uno por equipo & Facilita el marco Scrum y la resolución de impedimentos. & No tiene autoridad sobre alcance ni fechas. \\
\end{longtable}

Las responsabilidades ágiles de dueño de producto y Scrum Master complementan esta estructura, pero no otorgan autoridad para aprobar cambios contractuales, aceptar hitos ni autorizar pasos a producción.

\subsubsection{Instancias de gobierno}\label{t9-2-instancias}

Las instancias funcionan con una cadencia estable y no crean niveles paralelos de aprobación. La Tabla~\mbox{\ref{tab:t9-instancias}} detalla quiénes las integran, su cadencia, su ámbito y el registro que dejan.

\begin{longtable}[]{@{}
  >{\raggedright\arraybackslash}p{(\columnwidth - 10\tabcolsep) * \real{0.1700}}
  >{\raggedright\arraybackslash}p{(\columnwidth - 10\tabcolsep) * \real{0.2500}}
  >{\raggedright\arraybackslash}p{(\columnwidth - 10\tabcolsep) * \real{0.1400}}
  >{\raggedright\arraybackslash}p{(\columnwidth - 10\tabcolsep) * \real{0.2700}}
  >{\raggedright\arraybackslash}p{(\columnwidth - 10\tabcolsep) * \real{0.1700}}@{}}
\caption{Instancias de gobierno del proyecto.}\label{tab:t9-instancias}\\
\toprule\noalign{}
\begin{minipage}[b]{\linewidth}\raggedright
\textbf{Instancia}
\end{minipage} & \begin{minipage}[b]{\linewidth}\raggedright
\textbf{Integrantes}
\end{minipage} & \begin{minipage}[b]{\linewidth}\raggedright
\textbf{Cadencia}
\end{minipage} & \begin{minipage}[b]{\linewidth}\raggedright
\textbf{Ámbito de decisión}
\end{minipage} & \begin{minipage}[b]{\linewidth}\raggedright
\textbf{Registro}
\end{minipage} \\
\midrule\noalign{}
\endfirsthead
\toprule\noalign{}
\begin{minipage}[b]{\linewidth}\raggedright
\textbf{Instancia}
\end{minipage} & \begin{minipage}[b]{\linewidth}\raggedright
\textbf{Integrantes}
\end{minipage} & \begin{minipage}[b]{\linewidth}\raggedright
\textbf{Cadencia}
\end{minipage} & \begin{minipage}[b]{\linewidth}\raggedright
\textbf{Ámbito de decisión}
\end{minipage} & \begin{minipage}[b]{\linewidth}\raggedright
\textbf{Registro}
\end{minipage} \\
\midrule\noalign{}
\endhead
\bottomrule\noalign{}
\multicolumn{5}{@{}l}{\rule{0pt}{3.2ex}\textbf{Fuente:} Elaboración propia.} \\
\endlastfoot
\textbf{Comité Ejecutivo} & Patrocinador, gerencia de Only Simple Solutions y administrador del contrato; el Jefe de Proyecto presenta. & Mensual & Decisiones estratégicas, asuntos escalados, cambios de alcance y riesgos mayores. Aprueba o rechaza las solicitudes de cambio. & Acta en SharePoint dentro de dos días hábiles. \\ \addlinespace
\textbf{Comité de Proyecto} & Contraparte técnica y Jefe de Proyecto; dueños de producto según la agenda. & Quincenal & Cronograma, entregables, desviaciones, riesgos, dependencias y decisiones operativas. Recomienda cambios y revisa la preparación de cada paso de fase. & Acta en SharePoint dentro de dos días hábiles. \\ \addlinespace
\textbf{Comité de Arquitectura} & Arquitecto de Solución, Encargado de Seguridad y referentes técnicos de Ancoa. & Mensual & Decisiones de arquitectura, deuda técnica, riesgos y controversias técnicas. & Acta y registro de ADR. \\ \addlinespace
\textbf{Comité de Operación} & Líder de Operación / SRE, mesa de servicio y contraparte técnica. & Mensual desde el mes 13 & Niveles de servicio, incidentes, problemas, continuidad y mejora continua. & Acta en SharePoint dentro de dos días hábiles. \\ \addlinespace
\textbf{Reunión semanal de seguimiento} & Equipos de ambas partes y responsables de frente. & Semanal & Impedimentos, dependencias y compromisos de ejecución. & Compromisos en Jira. \\ \addlinespace
\textbf{Ceremonias Scrum} & Equipo Scrum; usuarios clave en la revisión del sprint. & Ciclo de dos semanas & Planificación, coordinación, revisión y mejora del trabajo del sprint, sin efecto contractual. & Jira. \\
\end{longtable}

Cada convocatoria de comité incluirá agenda y datos. La Tabla~\mbox{\ref{tab:t9-agendas}} fija las entradas y la agenda mínima de los cuatro comités, para que cada sesión llegue con la información que necesita para decidir.

\begin{longtable}[]{@{}
  >{\raggedright\arraybackslash}p{(\columnwidth - 6\tabcolsep) * \real{0.2000}}
  >{\raggedright\arraybackslash}p{(\columnwidth - 6\tabcolsep) * \real{0.3800}}
  >{\raggedright\arraybackslash}p{(\columnwidth - 6\tabcolsep) * \real{0.4200}}@{}}
\caption{Entradas y agenda mínima de los comités.}\label{tab:t9-agendas}\\
\toprule\noalign{}
\begin{minipage}[b]{\linewidth}\raggedright
\textbf{Comité}
\end{minipage} & \begin{minipage}[b]{\linewidth}\raggedright
\textbf{Entradas}
\end{minipage} & \begin{minipage}[b]{\linewidth}\raggedright
\textbf{Agenda mínima}
\end{minipage} \\
\midrule\noalign{}
\endfirsthead
\toprule\noalign{}
\begin{minipage}[b]{\linewidth}\raggedright
\textbf{Comité}
\end{minipage} & \begin{minipage}[b]{\linewidth}\raggedright
\textbf{Entradas}
\end{minipage} & \begin{minipage}[b]{\linewidth}\raggedright
\textbf{Agenda mínima}
\end{minipage} \\
\midrule\noalign{}
\endhead
\bottomrule\noalign{}
\multicolumn{3}{@{}l}{\rule{0pt}{3.2ex}\textbf{Fuente:} Elaboración propia.} \\
\endlastfoot
\textbf{Comité Ejecutivo} & Informe mensual de desempeño, solicitudes de cambio recomendadas, riesgos mayores y asuntos escalados. & Avance contra hitos; valor ganado; decisiones de cambio; riesgos mayores; asuntos escalados; acuerdos y responsables. \\ \addlinespace
\textbf{Comité de Proyecto} & Tablero de estado, registro de riesgos, registro de dependencias, solicitudes de cambio y estado del involucramiento. & Cronograma y entregables; desviaciones y planes de recuperación; riesgos; dependencias; análisis de cambios; preparación de fase. \\ \addlinespace
\textbf{Comité de Arquitectura} & ADR propuestas, registro de deuda técnica, resultados de funciones de aptitud y dependencias técnicas. & Aprobación de ADR; estado de la deuda; riesgos técnicos; cambios informados por fabricantes y proveedores. \\ \addlinespace
\textbf{Comité de Operación} & Informes de nivel de servicio, registro de incidentes y problemas, capacidad y resultados de pruebas de continuidad. & Cumplimiento de niveles de servicio; incidentes y problemas; continuidad; mejora continua; preparación de la marcha blanca. \\
\end{longtable}

\subsubsection{Escalamiento y reglas de decisión}\label{t9-2-escalamiento}

Cada materia tiene una ruta de escalamiento definida. La Tabla~\mbox{\ref{tab:t9-escalamiento}} indica el primer nivel que la trata, el nivel siguiente, la instancia final y el plazo o la regla que se aplica.

\begin{longtable}[]{@{}
  >{\raggedright\arraybackslash}p{(\columnwidth - 10\tabcolsep) * \real{0.2000}}
  >{\raggedright\arraybackslash}p{(\columnwidth - 10\tabcolsep) * \real{0.2000}}
  >{\raggedright\arraybackslash}p{(\columnwidth - 10\tabcolsep) * \real{0.1800}}
  >{\raggedright\arraybackslash}p{(\columnwidth - 10\tabcolsep) * \real{0.2000}}
  >{\raggedright\arraybackslash}p{(\columnwidth - 10\tabcolsep) * \real{0.2200}}@{}}
\caption{Rutas de escalamiento por materia.}\label{tab:t9-escalamiento}\\
\toprule\noalign{}
\begin{minipage}[b]{\linewidth}\raggedright
\textbf{Materia}
\end{minipage} & \begin{minipage}[b]{\linewidth}\raggedright
\textbf{Primer nivel}
\end{minipage} & \begin{minipage}[b]{\linewidth}\raggedright
\textbf{Segundo nivel}
\end{minipage} & \begin{minipage}[b]{\linewidth}\raggedright
\textbf{Instancia final}
\end{minipage} & \begin{minipage}[b]{\linewidth}\raggedright
\textbf{Plazo o regla}
\end{minipage} \\
\midrule\noalign{}
\endfirsthead
\toprule\noalign{}
\begin{minipage}[b]{\linewidth}\raggedright
\textbf{Materia}
\end{minipage} & \begin{minipage}[b]{\linewidth}\raggedright
\textbf{Primer nivel}
\end{minipage} & \begin{minipage}[b]{\linewidth}\raggedright
\textbf{Segundo nivel}
\end{minipage} & \begin{minipage}[b]{\linewidth}\raggedright
\textbf{Instancia final}
\end{minipage} & \begin{minipage}[b]{\linewidth}\raggedright
\textbf{Plazo o regla}
\end{minipage} \\
\midrule\noalign{}
\endhead
\bottomrule\noalign{}
\multicolumn{5}{@{}l}{\rule{0pt}{3.2ex}\textbf{Fuente:} Elaboración propia.} \\
\endlastfoot
\textbf{Bloqueo operativo o dependencia} & Responsable del frente, en la reunión semanal & Comité de Proyecto & Comité Ejecutivo & Sesión siguiente; si compromete un hito, aviso inmediato al Jefe de Proyecto. \\ \addlinespace
\textbf{Controversia técnica} & Arquitecto de Solución con los responsables técnicos & Comité de Arquitectura & Comité Ejecutivo & Sesión siguiente del Comité de Arquitectura. \\ \addlinespace
\textbf{Conflicto entre interesados} & Dueño de producto con el Jefe de Proyecto & Comité de Proyecto & Comité Ejecutivo, si es entre gerencias o afecta un hito & Cinco días hábiles en el primer nivel; luego, sesión siguiente. \\ \addlinespace
\textbf{Asunto contractual o fuera de atribuciones} & Jefe de Proyecto & -- & Comité Ejecutivo & Acuerdo escrito de ambas partes cuando modifica el contrato. \\ \addlinespace
\textbf{Cruce de datos Retail--Emisor} & Líder de Datos, Encargado de Seguridad y contraparte de cumplimiento & Comité de Arquitectura & Comité Ejecutivo & No se aprueba sin las tres evaluaciones documentadas. \\ \addlinespace
\textbf{Amenaza a seguridad, integridad de datos o continuidad} & Encargado de Seguridad o Líder de Operación & Jefe de Proyecto & Contrapartes autorizadas de Ancoa & Comunicación inmediata; incidente crítico en un máximo de dos horas. \\
\end{longtable}

Las decisiones de cualquier instancia siguen cinco reglas.

\begin{enumerate}
\tightlist
\item
  No existe aprobación por silencio. Si falta un responsable necesario o hay desacuerdo, el asunto queda abierto y se escala.
\item
  Las decisiones conjuntas requieren representantes autorizados de ambas partes y los responsables de la materia afectada.
\item
  El acuerdo es explícito y queda en acta, emitida dentro de los dos días hábiles siguientes, con responsables y plazos.
\item
  Un asunto escalado sigue abierto hasta comprobar el acuerdo y su ejecución.
\item
  Ninguna prioridad comercial autoriza eliminar la separación entre el negocio retail y el emisor.
\end{enumerate}

\subsubsection{Matriz de decisiones de gestión}\label{t9-2-matriz}

La Tabla~\mbox{\ref{tab:t9-decisiones}} identifica, para cada decisión de gestión, quién la prepara, quién la revisa, quién la aprueba y dónde queda registrada. Las decisiones del desarrollo se detallan en el Formulario T-10.

\begin{longtable}[]{@{}
  >{\raggedright\arraybackslash}p{(\columnwidth - 10\tabcolsep) * \real{0.2200}}
  >{\raggedright\arraybackslash}p{(\columnwidth - 10\tabcolsep) * \real{0.2000}}
  >{\raggedright\arraybackslash}p{(\columnwidth - 10\tabcolsep) * \real{0.2000}}
  >{\raggedright\arraybackslash}p{(\columnwidth - 10\tabcolsep) * \real{0.2000}}
  >{\raggedright\arraybackslash}p{(\columnwidth - 10\tabcolsep) * \real{0.1800}}@{}}
\caption{Matriz de decisiones de gestión.}\label{tab:t9-decisiones}\\
\toprule\noalign{}
\begin{minipage}[b]{\linewidth}\raggedright
\textbf{Decisión}
\end{minipage} & \begin{minipage}[b]{\linewidth}\raggedright
\textbf{Prepara}
\end{minipage} & \begin{minipage}[b]{\linewidth}\raggedright
\textbf{Revisa}
\end{minipage} & \begin{minipage}[b]{\linewidth}\raggedright
\textbf{Aprueba}
\end{minipage} & \begin{minipage}[b]{\linewidth}\raggedright
\textbf{Registro}
\end{minipage} \\
\midrule\noalign{}
\endfirsthead
\toprule\noalign{}
\begin{minipage}[b]{\linewidth}\raggedright
\textbf{Decisión}
\end{minipage} & \begin{minipage}[b]{\linewidth}\raggedright
\textbf{Prepara}
\end{minipage} & \begin{minipage}[b]{\linewidth}\raggedright
\textbf{Revisa}
\end{minipage} & \begin{minipage}[b]{\linewidth}\raggedright
\textbf{Aprueba}
\end{minipage} & \begin{minipage}[b]{\linewidth}\raggedright
\textbf{Registro}
\end{minipage} \\
\midrule\noalign{}
\endhead
\bottomrule\noalign{}
\multicolumn{5}{@{}l}{\rule{0pt}{3.2ex}\textbf{Fuente:} Elaboración propia.} \\
\endlastfoot
\textbf{Plan para la dirección y líneas base (H1)} & Jefe de Proyecto & Líder Funcional, contrapartes usuarias y líderes técnicos & Comité Ejecutivo & Versión aprobada en SharePoint. \\ \addlinespace
\textbf{Solicitud de cambio} & Solicitante; análisis de impacto del Jefe de Proyecto & Comité de Proyecto (recomendación) & Comité Ejecutivo; ambas partes por escrito si es contractual & Registro de cambios. \\ \addlinespace
\textbf{Respuesta a un riesgo que afecta una línea base} & Dueño del riesgo & Comité de Proyecto & Comité Ejecutivo, por control integrado de cambios & Registros de riesgos y de cambios. \\ \addlinespace
\textbf{Plan de recuperación por desviación superior al 10 \%} & Jefe de Proyecto & Comité de Proyecto & Comité Ejecutivo cuando excede las atribuciones del Jefe de Proyecto & Aviso formal y plan. \\ \addlinespace
\textbf{Aceptación de un hito} & Jefe de Proyecto, con entregable, evidencia y trazabilidad & Comité de Proyecto & Contraparte técnica & Acta de aceptación. \\ \addlinespace
\textbf{Cierre de marcha blanca y paso a producción} & Jefe de Proyecto, con las evaluaciones de calidad, seguridad, datos, implantación y operación & Comité de Proyecto & Contraparte técnica & Acta de aceptación. \\ \addlinespace
\textbf{Orden del backlog} & Dueño de producto & Equipo y Líder de Desarrollo & Dueño de producto, dentro del alcance aprobado & Jira. \\ \addlinespace
\textbf{Modificación contractual} & Administrador del contrato y Jefe de Proyecto & Comité Ejecutivo & Ambas partes, por escrito y dentro del límite acumulado del 20 \% & Registro de cambios. \\
\end{longtable}

\subsection{T-9.3 Integración y control de cambios}\label{t9-3}

\subsubsection{Procesos de gestión de la integración}\label{t9-3-procesos}

Only Simple Solutions aplicará los siete procesos de gestión de la integración (PMI, 2017). La Tabla~\mbox{\ref{tab:t9-integracion}} indica cuándo se ejecuta cada uno, quién responde y qué salida deja.

\begin{longtable}[]{@{}
  >{\raggedright\arraybackslash}p{(\columnwidth - 8\tabcolsep) * \real{0.2400}}
  >{\raggedright\arraybackslash}p{(\columnwidth - 8\tabcolsep) * \real{0.2400}}
  >{\raggedright\arraybackslash}p{(\columnwidth - 8\tabcolsep) * \real{0.2200}}
  >{\raggedright\arraybackslash}p{(\columnwidth - 8\tabcolsep) * \real{0.3000}}@{}}
\caption{Procesos de gestión de la integración.}\label{tab:t9-integracion}\\
\toprule\noalign{}
\begin{minipage}[b]{\linewidth}\raggedright
\textbf{Proceso}
\end{minipage} & \begin{minipage}[b]{\linewidth}\raggedright
\textbf{Momento}
\end{minipage} & \begin{minipage}[b]{\linewidth}\raggedright
\textbf{Responsable}
\end{minipage} & \begin{minipage}[b]{\linewidth}\raggedright
\textbf{Salida y evidencia}
\end{minipage} \\
\midrule\noalign{}
\endfirsthead
\toprule\noalign{}
\begin{minipage}[b]{\linewidth}\raggedright
\textbf{Proceso}
\end{minipage} & \begin{minipage}[b]{\linewidth}\raggedright
\textbf{Momento}
\end{minipage} & \begin{minipage}[b]{\linewidth}\raggedright
\textbf{Responsable}
\end{minipage} & \begin{minipage}[b]{\linewidth}\raggedright
\textbf{Salida y evidencia}
\end{minipage} \\
\midrule\noalign{}
\endhead
\bottomrule\noalign{}
\multicolumn{4}{@{}l}{\rule{0pt}{3.2ex}\textbf{Fuente:} Elaboración propia a partir de PMI (2017).} \\
\endlastfoot
\textbf{Desarrollar el acta de constitución} & Inicio del contrato, mes 1 & Jefe de Proyecto con el patrocinador y los representantes autorizados & Acta de constitución; registro de supuestos, restricciones y dependencias. \\ \addlinespace
\textbf{Desarrollar el plan para la dirección del proyecto} & Mes 1; aprobación en H1 & Jefe de Proyecto con los responsables técnicos y funcionales & Plan integrado con sus planes subsidiarios y las líneas base de alcance, cronograma y costos. \\ \addlinespace
\textbf{Dirigir y gestionar el trabajo} & Toda la implementación & Jefe de Proyecto; equipos Scrum y responsables de flujos Kanban & Entregables, datos de desempeño y registros de incidencias e impedimentos. \\ \addlinespace
\textbf{Gestionar el conocimiento} & Cierre de sprint, de intervalo, de ensayo y de ola & Responsable de cada frente & Registro de lecciones aprendidas, catálogo de ADR y sesiones de transferencia. \\ \addlinespace
\textbf{Monitorear y controlar el trabajo} & Semanal, quincenal y mensual & Jefe de Proyecto & Informe mensual con valor ganado; acciones correctivas, preventivas y reparación de defectos. \\ \addlinespace
\textbf{Realizar el control integrado de cambios} & Ante cada solicitud & Jefe de Proyecto; Comité de Proyecto y Comité Ejecutivo & Registro de cambios y líneas base actualizadas, con sus versiones anteriores. \\ \addlinespace
\textbf{Cerrar el proyecto o fase} & Mes 16 (Etapa 1), mes 21 (implementación) y mes 56 (contrato) & Jefe de Proyecto; la contraparte técnica formaliza & Actas de aceptación, lecciones aprendidas y transferencia al frente que continúa. \\
\end{longtable}

El acta de constitución no sustituye el contrato ni permite omitir obligaciones. Las versiones aprobadas del plan se conservan en SharePoint con fecha, responsable y decisión de aprobación. El avance y el pronóstico se actualizan sin reemplazar las líneas base, para mantener visibles las desviaciones.

\subsubsection{Procedimiento de control integrado de cambios}\label{t9-3-cambios}

Una necesidad funcional, un hallazgo técnico, una incidencia o una observación del cliente puede originar una solicitud, pero su registro no autoriza la ejecución. La Tabla~\mbox{\ref{tab:t9-cambios}} describe los cinco pasos del procedimiento.

\begin{longtable}[]{@{}
  >{\raggedright\arraybackslash}p{(\columnwidth - 8\tabcolsep) * \real{0.1600}}
  >{\raggedright\arraybackslash}p{(\columnwidth - 8\tabcolsep) * \real{0.2200}}
  >{\raggedright\arraybackslash}p{(\columnwidth - 8\tabcolsep) * \real{0.4200}}
  >{\raggedright\arraybackslash}p{(\columnwidth - 8\tabcolsep) * \real{0.2000}}@{}}
\caption{Procedimiento de control integrado de cambios.}\label{tab:t9-cambios}\\
\toprule\noalign{}
\begin{minipage}[b]{\linewidth}\raggedright
\textbf{Paso}
\end{minipage} & \begin{minipage}[b]{\linewidth}\raggedright
\textbf{Responsable}
\end{minipage} & \begin{minipage}[b]{\linewidth}\raggedright
\textbf{Actividades}
\end{minipage} & \begin{minipage}[b]{\linewidth}\raggedright
\textbf{Registro}
\end{minipage} \\
\midrule\noalign{}
\endfirsthead
\toprule\noalign{}
\begin{minipage}[b]{\linewidth}\raggedright
\textbf{Paso}
\end{minipage} & \begin{minipage}[b]{\linewidth}\raggedright
\textbf{Responsable}
\end{minipage} & \begin{minipage}[b]{\linewidth}\raggedright
\textbf{Actividades}
\end{minipage} & \begin{minipage}[b]{\linewidth}\raggedright
\textbf{Registro}
\end{minipage} \\
\midrule\noalign{}
\endhead
\bottomrule\noalign{}
\multicolumn{4}{@{}l}{\rule{0pt}{3.2ex}\textbf{Fuente:} Elaboración propia.} \\
\endlastfoot
\textbf{1. Registro} & Solicitante; el Jefe de Proyecto asigna el identificador & Describir el motivo, el resultado esperado y los entregables afectados. Registrar fecha, origen y antecedentes. Vincular la solicitud con requisitos, observaciones, riesgos o incidencias. & Solicitud en estado registrada. \\ \addlinespace
\textbf{2. Análisis de impacto} & Jefe de Proyecto y líderes correspondientes & Evaluar alcance, plazo, costo, riesgo y niveles de servicio, además de arquitectura, interfaces, datos, seguridad, calidad, recursos, capacitación y operación. Identificar dependencias, alternativas y consecuencias de aprobar o rechazar. Si involucra datos u obligaciones del emisor, incorporar las evaluaciones del Líder de Datos, del Encargado de Seguridad y de la contraparte de cumplimiento. & Informe de impacto. \\ \addlinespace
\textbf{3. Decisión} & Comité de Proyecto recomienda; Comité Ejecutivo aprueba o rechaza & Decidir antes de la ejecución. Si existe modificación contractual, obtener acuerdo escrito de ambas partes. & Acta y estado de la solicitud. \\ \addlinespace
\textbf{4. Actualización y comunicación} & Jefe de Proyecto y responsables de cada registro & Actualizar líneas base y documentos, conservando las versiones anteriores y el fundamento. Ajustar requisitos, ADR, backlogs, procedimientos y pruebas. Comunicar qué cambia, desde cuándo y qué acciones deben realizarse. & Nueva versión en SharePoint. \\ \addlinespace
\textbf{5. Ejecución y cierre} & Responsable designado; el Jefe de Proyecto comprueba & Ejecutar el cambio y aportar evidencia de su resultado. La aceptación del entregable afectado corresponde a la contraparte técnica. Una solicitud rechazada conserva su justificación y se comunica al solicitante. & Solicitud cerrada. \\
\end{longtable}

El procedimiento se aplica con las restricciones siguientes.

\begin{itemize}
\tightlist
\item
  No se admiten modificaciones que alteren el cronograma obligatorio (producción de la Etapa 1 desde el mes 16, producción de la Etapa 2 e inicio de la operación en el mes 21 y término en el mes 56), el objeto principal, la naturaleza del contrato o las garantías mínimas.
\item
  Las modificaciones contractuales respetan el límite acumulado del 20 \% del valor original del contrato. El registro identifica el efecto acumulado y la autorización, sin publicar importes en la Oferta Técnica.
\item
  El refinamiento del backlog que divide tareas, aclara detalles u ordena la construcción, manteniendo requisitos, entregables, controles y fechas, no requiere solicitud. Si afecta una línea base o sus consecuencias son inciertas, se somete a evaluación formal.
\item
  La atención urgente de una incidencia usa los procedimientos de respuesta y reversión autorizados. La urgencia no se convierte en aprobación de nuevos compromisos.
\end{itemize}

\subsubsection{Medición del desempeño con valor ganado}\label{t9-3-evm}

El valor ganado relaciona el trabajo verificado de la EDT con su presupuesto y su calendario de control (Prasad et al., 2006). La Tabla~\mbox{\ref{tab:t9-evm}} define los indicadores y la regla con que se usan.

\begin{longtable}[]{@{}
  >{\raggedright\arraybackslash}p{(\columnwidth - 8\tabcolsep) * \real{0.2200}}
  >{\raggedright\arraybackslash}p{(\columnwidth - 8\tabcolsep) * \real{0.3400}}
  >{\raggedright\arraybackslash}p{(\columnwidth - 8\tabcolsep) * \real{0.3000}}
  >{\raggedright\arraybackslash}p{(\columnwidth - 8\tabcolsep) * \real{0.1400}}@{}}
\caption{Indicadores de valor ganado.}\label{tab:t9-evm}\\
\toprule\noalign{}
\begin{minipage}[b]{\linewidth}\raggedright
\textbf{Indicador}
\end{minipage} & \begin{minipage}[b]{\linewidth}\raggedright
\textbf{Definición o fórmula}
\end{minipage} & \begin{minipage}[b]{\linewidth}\raggedright
\textbf{Regla de uso}
\end{minipage} & \begin{minipage}[b]{\linewidth}\raggedright
\textbf{Frecuencia}
\end{minipage} \\
\midrule\noalign{}
\endfirsthead
\toprule\noalign{}
\begin{minipage}[b]{\linewidth}\raggedright
\textbf{Indicador}
\end{minipage} & \begin{minipage}[b]{\linewidth}\raggedright
\textbf{Definición o fórmula}
\end{minipage} & \begin{minipage}[b]{\linewidth}\raggedright
\textbf{Regla de uso}
\end{minipage} & \begin{minipage}[b]{\linewidth}\raggedright
\textbf{Frecuencia}
\end{minipage} \\
\midrule\noalign{}
\endhead
\bottomrule\noalign{}
\multicolumn{4}{@{}l}{\rule{0pt}{3.2ex}\textbf{Fuente:} Elaboración propia a partir de PMI (2017) y Prasad et al. (2006).} \\
\endlastfoot
\textbf{Valor planificado (PV)} & Valor presupuestado del trabajo previsto hasta la fecha de corte. & Según la línea base de costos y el calendario de control. & Mensual \\ \addlinespace
\textbf{Valor ganado (EV)} & Valor presupuestado del trabajo efectivamente realizado y comprobado. & Se reconoce solo con evidencia, según la regla del paquete. & Mensual \\ \addlinespace
\textbf{Costo real (AC)} & Costo incurrido por el trabajo realizado. & A partir de los antecedentes de costos autorizados. & Mensual \\ \addlinespace
\textbf{Índice de desempeño del cronograma (SPI)} & SPI = EV / PV & Un valor menor que 1 indica atraso. Si PV es cero, se registra que no es calculable y su causa. & Mensual \\ \addlinespace
\textbf{Índice de desempeño del costo (CPI)} & CPI = EV / AC & Un valor menor que 1 indica sobrecosto. Si AC es cero, se registra que no es calculable y su causa. & Mensual \\
\end{longtable}

La medición sigue cuatro reglas, que evitan declarar avance con porcentajes subjetivos.

\begin{enumerate}
\tightlist
\item
  Antes de iniciar cada paquete de trabajo, el Jefe de Proyecto fija su regla de reconocimiento. Para paquetes de hasta un mes se usa la regla 0/100, en que el valor se reconoce al verificar el producto. Para paquetes más largos se usan hitos ponderados, cada uno con un producto verificable y un porcentaje definido de antemano.
\item
  El reconocimiento de valor ganado para control no equivale a la aceptación contractual ni al derecho a pago.
\item
  Las métricas de Scrum y Kanban complementan el seguimiento de capacidad, bloqueos y entrega, sin sustituir el valor ganado.
\item
  Toda desviación superior al 10 \% en un hito, respecto de su línea base de control, se comunica dentro de los cinco días hábiles siguientes a su detección, con un plan de recuperación.
\end{enumerate}

El informe mensual presenta los índices con su fecha de corte, su evidencia y la explicación de las desviaciones, con trazabilidad hacia los paquetes de trabajo. Los antecedentes financieros se comunican por los canales contractuales autorizados.

\subsubsection{Condiciones de cierre de marcha blanca y de fase}\label{t9-3-cierre}

El paso a producción de cada etapa exige que, al cierre de su marcha blanca, se cumplan de forma copulativa las cinco condiciones de la Tabla~\mbox{\ref{tab:t9-cierre}}. Una revisión de sprint o una declaración interna de término no puede sustituirlas.

\begin{longtable}[]{@{}
  >{\raggedright\arraybackslash}p{(\columnwidth - 6\tabcolsep) * \real{0.4600}}
  >{\raggedright\arraybackslash}p{(\columnwidth - 6\tabcolsep) * \real{0.3000}}
  >{\raggedright\arraybackslash}p{(\columnwidth - 6\tabcolsep) * \real{0.2400}}@{}}
\caption{Condiciones copulativas de cierre de marcha blanca.}\label{tab:t9-cierre}\\
\toprule\noalign{}
\begin{minipage}[b]{\linewidth}\raggedright
\textbf{Condición}
\end{minipage} & \begin{minipage}[b]{\linewidth}\raggedright
\textbf{Verificación}
\end{minipage} & \begin{minipage}[b]{\linewidth}\raggedright
\textbf{Responsable de la evidencia}
\end{minipage} \\
\midrule\noalign{}
\endfirsthead
\toprule\noalign{}
\begin{minipage}[b]{\linewidth}\raggedright
\textbf{Condición}
\end{minipage} & \begin{minipage}[b]{\linewidth}\raggedright
\textbf{Verificación}
\end{minipage} & \begin{minipage}[b]{\linewidth}\raggedright
\textbf{Responsable de la evidencia}
\end{minipage} \\
\midrule\noalign{}
\endhead
\bottomrule\noalign{}
\multicolumn{3}{@{}l}{\rule{0pt}{3.2ex}\textbf{Fuente:} Elaboración propia.} \\
\endlastfoot
1. No hay incidentes críticos ni altos abiertos atribuibles a la solución. & Registro de incidentes. & Líder de Operación / SRE \\ \addlinespace
2. El volumen real comprometido y los indicadores de disponibilidad y respuesta se sostuvieron durante al menos las últimas cuatro semanas. & Plataforma de observabilidad e informes de servicio. & Líder de Operación / SRE \\ \addlinespace
3. La conciliación no presenta diferencias no explicadas. & Informes de conciliación. & Líder de Datos \\ \addlinespace
4. El personal del cliente recibió la capacitación y la certificación exigidas. & Registros de la plataforma de capacitación. & Líder de Implantación y Gestión del Cambio \\ \addlinespace
5. La contraparte técnica suscribió el acta de aceptación. & Acta de aceptación. & Contraparte técnica \\
\end{longtable}

Si al cierre del período no se cumplen estas condiciones, Only Simple Solutions aplicará la recuperación y extensión de marcha blanca que exige el contrato, a su cargo, sin desplazar las fases siguientes ni dispensar las responsabilidades por atraso. Para cada acta de aceptación, el cliente dispondrá de diez días hábiles para revisar y pronunciarse, y Only Simple Solutions de diez días hábiles para subsanar las observaciones.

El cierre de implementación del mes 21 transfiere ambos alcances a operación, con sus procedimientos, documentación y responsabilidades. El cierre contractual del mes 56 verifica las obligaciones finales de transferencia y reversibilidad. Ningún cierre extingue responsabilidades que deban mantenerse en el período siguiente.

\subsection{T-9.4 Gestión de interesados}\label{t9-4}

\subsubsection{Registro de interesados}\label{t9-4-registro}

La Tabla~\mbox{\ref{tab:t9-registro}} presenta el universo analítico de 26 interesados con su grupo, su clasificación según los criterios de la sección 6.1.4.1 y su cuadrante en la matriz de poder e interés. La clasificación indica, en este orden, la relación con Ancoa (interno o externo), la forma de participación (actor o no operativo) y el contexto.

\begin{longtable}[]{@{}
  >{\raggedright\arraybackslash}p{(\columnwidth - 10\tabcolsep) * \real{0.0800}}
  >{\raggedright\arraybackslash}p{(\columnwidth - 10\tabcolsep) * \real{0.2800}}
  >{\raggedright\arraybackslash}p{(\columnwidth - 10\tabcolsep) * \real{0.2200}}
  >{\raggedright\arraybackslash}p{(\columnwidth - 10\tabcolsep) * \real{0.2400}}
  >{\raggedright\arraybackslash}p{(\columnwidth - 10\tabcolsep) * \real{0.1800}}@{}}
\caption{Registro de interesados: identificación y clasificación.}\label{tab:t9-registro}\\
\toprule\noalign{}
\begin{minipage}[b]{\linewidth}\raggedright
\textbf{ID}
\end{minipage} & \begin{minipage}[b]{\linewidth}\raggedright
\textbf{Interesado}
\end{minipage} & \begin{minipage}[b]{\linewidth}\raggedright
\textbf{Grupo de interés}
\end{minipage} & \begin{minipage}[b]{\linewidth}\raggedright
\textbf{Clasificación}
\end{minipage} & \begin{minipage}[b]{\linewidth}\raggedright
\textbf{Cuadrante}
\end{minipage} \\
\midrule\noalign{}
\endfirsthead
\toprule\noalign{}
\begin{minipage}[b]{\linewidth}\raggedright
\textbf{ID}
\end{minipage} & \begin{minipage}[b]{\linewidth}\raggedright
\textbf{Interesado}
\end{minipage} & \begin{minipage}[b]{\linewidth}\raggedright
\textbf{Grupo de interés}
\end{minipage} & \begin{minipage}[b]{\linewidth}\raggedright
\textbf{Clasificación}
\end{minipage} & \begin{minipage}[b]{\linewidth}\raggedright
\textbf{Cuadrante}
\end{minipage} \\
\midrule\noalign{}
\endhead
\bottomrule\noalign{}
\multicolumn{5}{@{}l}{\rule{0pt}{3.2ex}\textbf{Fuente:} Elaboración propia a partir de las Figuras 6.1 y 6.2 del Subdocumento 6.} \\
\endlastfoot
\textbf{I-01} & Directorio & Dirección y control & Interno · no operativo · frontera & Gestionar de cerca \\ \addlinespace
\textbf{I-02} & Gerencia General & Dirección y control & Interno · no operativo · frontera & Gestionar de cerca \\ \addlinespace
\textbf{I-03} & Contraloría & Dirección y control & Interno · actor · frontera & Gestionar de cerca \\ \addlinespace
\textbf{I-04} & Grupo familiar controlador & Propiedad & Externo · no operativo · propiedad & Mantener satisfechos \\ \addlinespace
\textbf{I-05} & Fondos de inversión & Propiedad & Externo · no operativo · propiedad & Mantener satisfechos \\ \addlinespace
\textbf{I-06} & Gerencia Comercial & Negocios y operación & Interno · actor · retail & Gestionar de cerca \\ \addlinespace
\textbf{I-07} & Gerencia de Canales Digitales & Negocios y operación & Interno · actor · retail & Gestionar de cerca \\ \addlinespace
\textbf{I-08} & Gerencia de Logística & Negocios y operación & Interno · actor · retail & Gestionar de cerca \\ \addlinespace
\textbf{I-09} & Negocio Financiero (filial emisora) & Negocios y operación & Interno · actor · emisor & Gestionar de cerca \\ \addlinespace
\textbf{I-10} & Tecnologías de la Información & Soporte y control operacional & Interno · actor · plataforma compartida & Gestionar de cerca \\ \addlinespace
\textbf{I-11} & Prevención de Pérdidas & Soporte y control operacional & Interno · actor · retail & Gestionar de cerca \\ \addlinespace
\textbf{I-12} & Gestión de Personas (coordinación de personas) & Soporte y control operacional & Interno · no operativo · plataforma compartida & Gestionar de cerca \\ \addlinespace
\textbf{I-13} & Jefaturas de tienda y departamento & Operación de tienda & Interno · actor · retail & Mantener informados \\ \addlinespace
\textbf{I-14} & Vendedores y cajeros & Operación de tienda & Interno · actor · retail & Mantener informados \\ \addlinespace
\textbf{I-15} & Personal de reposición y bodega & Operación de tienda & Interno · actor · retail & Mantener informados \\ \addlinespace
\textbf{I-16} & Clientes & Clientes y terceros operacionales & Externo · actor · retail & Mantener informados \\ \addlinespace
\textbf{I-17} & Titulares de tarjeta & Clientes y terceros operacionales & Externo · actor · emisor & Mantener informados \\ \addlinespace
\textbf{I-18} & Vendedores de marketplace & Clientes y terceros operacionales & Externo · actor · retail & Mantener informados \\ \addlinespace
\textbf{I-19} & Proveedores tecnológicos & Clientes y terceros operacionales & Externo · actor · plataforma compartida & Mantener informados \\ \addlinespace
\textbf{I-20} & Proveedores logísticos & Clientes y terceros operacionales & Externo · actor · retail & Mantener informados \\ \addlinespace
\textbf{I-21} & Repositores externos & Clientes y terceros operacionales & Externo · actor · retail & Monitorear \\ \addlinespace
\textbf{I-22} & Administradores de centros comerciales & Clientes y terceros operacionales & Externo · no operativo · retail & Monitorear \\ \addlinespace
\textbf{I-23} & Comisión para el Mercado Financiero (CMF) & Organismos externos & Externo · no operativo · emisor & Gestionar de cerca \\ \addlinespace
\textbf{I-24} & Servicio Nacional del Consumidor (SERNAC) & Organismos externos & Externo · no operativo · frontera & Gestionar de cerca \\ \addlinespace
\textbf{I-25} & Servicio de Impuestos Internos (SII) & Organismos externos & Externo · no operativo · retail & Mantener satisfechos \\ \addlinespace
\textbf{I-26} & Agencia de Protección de Datos Personales & Organismos externos & Externo · no operativo · frontera & Mantener satisfechos \\
\end{longtable}

El registro reúne 13 interesados internos y 13 externos, con Propiedad considerada externa a la estructura operativa. Los roles de gobierno de la sección T-9.2 corresponden a personas de estos interesados: el patrocinador pertenece a I-02, la contraparte técnica a I-10 y la contraparte de cumplimiento a I-03 e I-09.

El registro completo conserva los tres bloques del PMBOK (PMI, 2017). La Tabla~\mbox{\ref{tab:t9-campos}} fija los campos de cada bloque, que se completarán con los representantes designados durante el levantamiento del mes 1.

\begin{longtable}[]{@{}
  >{\raggedright\arraybackslash}p{(\columnwidth - 6\tabcolsep) * \real{0.2000}}
  >{\raggedright\arraybackslash}p{(\columnwidth - 6\tabcolsep) * \real{0.5200}}
  >{\raggedright\arraybackslash}p{(\columnwidth - 6\tabcolsep) * \real{0.2800}}@{}}
\caption{Campos del registro de interesados.}\label{tab:t9-campos}\\
\toprule\noalign{}
\begin{minipage}[b]{\linewidth}\raggedright
\textbf{Bloque}
\end{minipage} & \begin{minipage}[b]{\linewidth}\raggedright
\textbf{Campos}
\end{minipage} & \begin{minipage}[b]{\linewidth}\raggedright
\textbf{Acceso}
\end{minipage} \\
\midrule\noalign{}
\endfirsthead
\toprule\noalign{}
\begin{minipage}[b]{\linewidth}\raggedright
\textbf{Bloque}
\end{minipage} & \begin{minipage}[b]{\linewidth}\raggedright
\textbf{Campos}
\end{minipage} & \begin{minipage}[b]{\linewidth}\raggedright
\textbf{Acceso}
\end{minipage} \\
\midrule\noalign{}
\endhead
\bottomrule\noalign{}
\multicolumn{3}{@{}l}{\rule{0pt}{3.2ex}\textbf{Fuente:} Elaboración propia a partir de PMI (2017).} \\
\endlastfoot
\textbf{Identificación} & Representante titular y suplente, cargo, área, contacto, atribuciones confirmadas por Ancoa y canal de validación. & Equipos del proyecto con función asignada. \\ \addlinespace
\textbf{Evaluación} & Requisitos y expectativas, influencia, interés, dirección de la influencia (ascendente, descendente, lateral o hacia afuera) y actitud observada. & Restringido al Jefe de Proyecto, al Líder de Implantación y al patrocinador. \\ \addlinespace
\textbf{Clasificación} & Relación con Ancoa, forma de participación, contexto, cuadrante de poder e interés, nivel actual y nivel deseado. & Equipos del proyecto con función asignada. \\
\end{longtable}

\subsubsection{Matriz de evaluación del involucramiento}\label{t9-4-matriz}

La Tabla~\mbox{\ref{tab:t9-involucramiento}} extiende a los 26 interesados la matriz de la Tabla 6.3 del Subdocumento 6. El nivel actual se determinará con entrevistas, encuestas y talleres del levantamiento. Los niveles deseados son la propuesta de Only Simple Solutions y se acordarán con Ancoa al aprobar la estrategia de involucramiento en H1. Los niveles usados son desconocedor, reticente, neutral, de apoyo y líder.

\begin{longtable}[]{@{}
  >{\raggedright\arraybackslash}p{(\columnwidth - 10\tabcolsep) * \real{0.0800}}
  >{\raggedright\arraybackslash}p{(\columnwidth - 10\tabcolsep) * \real{0.2400}}
  >{\raggedright\arraybackslash}p{(\columnwidth - 10\tabcolsep) * \real{0.1600}}
  >{\raggedright\arraybackslash}p{(\columnwidth - 10\tabcolsep) * \real{0.1400}}
  >{\raggedright\arraybackslash}p{(\columnwidth - 10\tabcolsep) * \real{0.3800}}@{}}
\caption{Matriz de evaluación del involucramiento.}\label{tab:t9-involucramiento}\\
\toprule\noalign{}
\begin{minipage}[b]{\linewidth}\raggedright
\textbf{ID}
\end{minipage} & \begin{minipage}[b]{\linewidth}\raggedright
\textbf{Interesado}
\end{minipage} & \begin{minipage}[b]{\linewidth}\raggedright
\textbf{Nivel actual}
\end{minipage} & \begin{minipage}[b]{\linewidth}\raggedright
\textbf{Nivel deseado}
\end{minipage} & \begin{minipage}[b]{\linewidth}\raggedright
\textbf{Acción de involucramiento}
\end{minipage} \\
\midrule\noalign{}
\endfirsthead
\toprule\noalign{}
\begin{minipage}[b]{\linewidth}\raggedright
\textbf{ID}
\end{minipage} & \begin{minipage}[b]{\linewidth}\raggedright
\textbf{Interesado}
\end{minipage} & \begin{minipage}[b]{\linewidth}\raggedright
\textbf{Nivel actual}
\end{minipage} & \begin{minipage}[b]{\linewidth}\raggedright
\textbf{Nivel deseado}
\end{minipage} & \begin{minipage}[b]{\linewidth}\raggedright
\textbf{Acción de involucramiento}
\end{minipage} \\
\midrule\noalign{}
\endhead
\bottomrule\noalign{}
\multicolumn{5}{@{}l}{\rule{0pt}{3.2ex}\textbf{Fuente:} Elaboración propia a partir de PMI (2017).} \\
\endlastfoot
\textbf{I-01} & Directorio & Por determinar en el mes 1 & De apoyo & Resumen por hito y decisiones mayores a través de la Gerencia General. \\ \addlinespace
\textbf{I-02} & Gerencia General & Por determinar en el mes 1 & Líder & Resolver prioridades y escalamientos en el Comité Ejecutivo. \\ \addlinespace
\textbf{I-03} & Contraloría & Por determinar en el mes 1 & De apoyo & Validar los controles de cruce y la evidencia de la frontera. \\ \addlinespace
\textbf{I-04} & Grupo familiar controlador & Por determinar en el mes 1 & Neutral & Información resumida por hito, canalizada por el Directorio. \\ \addlinespace
\textbf{I-05} & Fondos de inversión & Por determinar en el mes 1 & Neutral & Información resumida por hito, canalizada por el Directorio. \\ \addlinespace
\textbf{I-06} & Gerencia Comercial & Por determinar en el mes 1 & De apoyo & Acordar reglas de precio y comisiones; ordenar el backlog de precio y catálogo. \\ \addlinespace
\textbf{I-07} & Gerencia de Canales Digitales & Por determinar en el mes 1 & De apoyo & Ordenar el backlog de pedidos y acordar la regla de publicación de existencia. \\ \addlinespace
\textbf{I-08} & Gerencia de Logística & Por determinar en el mes 1 & De apoyo & Ordenar el backlog de inventario y abastecimiento; validar la exactitud por ola. \\ \addlinespace
\textbf{I-09} & Negocio Financiero & Por determinar en el mes 1 & De apoyo & Ordenar el backlog de crédito y validar la secuencia de migración de cartera. \\ \addlinespace
\textbf{I-10} & Tecnologías de la Información & Por determinar en el mes 1 & Líder & Ejercer la contraparte técnica y participar en las sesiones de plataforma y datos. \\ \addlinespace
\textbf{I-11} & Prevención de Pérdidas & Por determinar en el mes 1 & De apoyo & Participar en las pruebas de inventario. \\ \addlinespace
\textbf{I-12} & Gestión de Personas & Por determinar en el mes 1 & De apoyo & Programar la formación por perfil y turno e incorporar el curso de inducción permanente. \\ \addlinespace
\textbf{I-13} & Jefaturas de tienda & Por determinar en el mes 1 & De apoyo & Participar en pilotos y adopción; designar al agente de cambio de su tienda. \\ \addlinespace
\textbf{I-14} & Vendedores y cajeros & Por determinar en el mes 1 & De apoyo & Probar en el piloto el flujo de venta con crédito; formación breve por turno. \\ \addlinespace
\textbf{I-15} & Reposición y bodega & Por determinar en el mes 1 & De apoyo & Formación práctica en conteo y reposición, con demostraciones guiadas. \\ \addlinespace
\textbf{I-16} & Clientes & Por determinar en el mes 1 & Neutral & Comunicación de los cambios visibles por los canales de Ancoa. \\ \addlinespace
\textbf{I-17} & Titulares de tarjeta & Por determinar en el mes 1 & Neutral & Plan de comunicación por ola de migración de cartera, validado por su claridad. \\ \addlinespace
\textbf{I-18} & Vendedores de marketplace & Por determinar en el mes 1 & De apoyo & Talleres de integración a través de Ancoa. \\ \addlinespace
\textbf{I-19} & Proveedores tecnológicos & Por determinar en el mes 1 & De apoyo & Verificar interfaces y actividades de transición. \\ \addlinespace
\textbf{I-20} & Proveedores logísticos & Por determinar en el mes 1 & De apoyo & Talleres de integración de los eventos de entrega a través de Ancoa. \\ \addlinespace
\textbf{I-21} & Repositores externos & Por determinar en el mes 1 & De apoyo & Capacitación a través de su empleador, con material de Only Simple Solutions. \\ \addlinespace
\textbf{I-22} & Administradores de centros comerciales & Por determinar en el mes 1 & Neutral & Aviso anticipado de obras, accesos y enlaces a través de Ancoa. \\ \addlinespace
\textbf{I-23} & CMF & No aplica & No aplica & Evidencia entregada por la filial emisora; control de requerimientos y plazos. \\ \addlinespace
\textbf{I-24} & SERNAC & No aplica & No aplica & Evidencia de precio exhibido y cobrado, entregada a través de Ancoa. \\ \addlinespace
\textbf{I-25} & SII & No aplica & No aplica & Evidencia de documentos tributarios, entregada a través de Ancoa. \\ \addlinespace
\textbf{I-26} & Agencia de Protección de Datos Personales & No aplica & No aplica & Evidencia de tratamiento de datos, entregada a través de Ancoa. \\
\end{longtable}

Las autoridades externas no se miden como personas que deban adoptar una actitud de apoyo al proveedor; para ellas se controlan requerimientos, antecedentes y plazos. Las entrevistas del levantamiento inicial anticipan que las jefaturas de tienda y los vendedores podrían partir como reticentes, por lo que concentrarán las primeras acciones. La matriz se actualizará al cierre de cada intervalo de coordinación.

\subsubsection{Dueños de producto, usuarios clave y agentes de cambio}\label{t9-4-po}

Ancoa designará un dueño de producto por dominio. La Tabla~\mbox{\ref{tab:t9-po}} indica la gerencia que lo designa y las validaciones en que participan sus usuarios clave.

\begin{longtable}[]{@{}
  >{\raggedright\arraybackslash}p{(\columnwidth - 6\tabcolsep) * \real{0.2400}}
  >{\raggedright\arraybackslash}p{(\columnwidth - 6\tabcolsep) * \real{0.2600}}
  >{\raggedright\arraybackslash}p{(\columnwidth - 6\tabcolsep) * \real{0.5000}}@{}}
\caption{Dueños de producto y validaciones por dominio.}\label{tab:t9-po}\\
\toprule\noalign{}
\begin{minipage}[b]{\linewidth}\raggedright
\textbf{Dominio}
\end{minipage} & \begin{minipage}[b]{\linewidth}\raggedright
\textbf{Gerencia que designa}
\end{minipage} & \begin{minipage}[b]{\linewidth}\raggedright
\textbf{Usuarios clave y validaciones}
\end{minipage} \\
\midrule\noalign{}
\endfirsthead
\toprule\noalign{}
\begin{minipage}[b]{\linewidth}\raggedright
\textbf{Dominio}
\end{minipage} & \begin{minipage}[b]{\linewidth}\raggedright
\textbf{Gerencia que designa}
\end{minipage} & \begin{minipage}[b]{\linewidth}\raggedright
\textbf{Usuarios clave y validaciones}
\end{minipage} \\
\midrule\noalign{}
\endhead
\bottomrule\noalign{}
\multicolumn{3}{@{}l}{\rule{0pt}{3.2ex}\textbf{Fuente:} Elaboración propia.} \\
\endlastfoot
\textbf{Precio y catálogo} & Gerencia Comercial & Usuarios clave comerciales en revisiones de sprint y pruebas de aceptación; coordinación del reetiquetado con las jefaturas. \\ \addlinespace
\textbf{Inventario y abastecimiento} & Gerencia de Logística & Usuarios clave de Logística; Prevención de Pérdidas participa en las pruebas de inventario. \\ \addlinespace
\textbf{Pedidos} & Gerencia de Canales Digitales & Usuarios clave digitales; talleres con vendedores de marketplace, proveedores y transportistas, a través de Ancoa. \\ \addlinespace
\textbf{Crédito} & Negocio Financiero & Usuarios clave de la filial; Contraloría y Negocio Financiero validan la separación de datos. \\
\end{longtable}

Antes del hito H6 habrá un agente de cambio en cada una de las 22 tiendas, elegido por su jefatura. Las jefaturas, vendedores y cajeros de las tiendas del piloto probarán la solución antes de cada ola.

\subsubsection{Estrategias de involucramiento por dimensión}\label{t9-4-estrategias}

La estrategia de involucramiento se organiza en cuatro dimensiones. La Tabla~\mbox{\ref{tab:t9-dimensiones}} detalla sus destinatarios, sus acciones y su responsable.

\begin{longtable}[]{@{}
  >{\raggedright\arraybackslash}p{(\columnwidth - 8\tabcolsep) * \real{0.1800}}
  >{\raggedright\arraybackslash}p{(\columnwidth - 8\tabcolsep) * \real{0.2200}}
  >{\raggedright\arraybackslash}p{(\columnwidth - 8\tabcolsep) * \real{0.4000}}
  >{\raggedright\arraybackslash}p{(\columnwidth - 8\tabcolsep) * \real{0.2000}}@{}}
\caption{Estrategias de involucramiento por dimensión.}\label{tab:t9-dimensiones}\\
\toprule\noalign{}
\begin{minipage}[b]{\linewidth}\raggedright
\textbf{Dimensión}
\end{minipage} & \begin{minipage}[b]{\linewidth}\raggedright
\textbf{Destinatarios}
\end{minipage} & \begin{minipage}[b]{\linewidth}\raggedright
\textbf{Acciones}
\end{minipage} & \begin{minipage}[b]{\linewidth}\raggedright
\textbf{Responsable}
\end{minipage} \\
\midrule\noalign{}
\endfirsthead
\toprule\noalign{}
\begin{minipage}[b]{\linewidth}\raggedright
\textbf{Dimensión}
\end{minipage} & \begin{minipage}[b]{\linewidth}\raggedright
\textbf{Destinatarios}
\end{minipage} & \begin{minipage}[b]{\linewidth}\raggedright
\textbf{Acciones}
\end{minipage} & \begin{minipage}[b]{\linewidth}\raggedright
\textbf{Responsable}
\end{minipage} \\
\midrule\noalign{}
\endhead
\bottomrule\noalign{}
\multicolumn{4}{@{}l}{\rule{0pt}{3.2ex}\textbf{Fuente:} Elaboración propia.} \\
\endlastfoot
\textbf{Comunicación} & Los siete grupos de interés & Comité Ejecutivo mensual para la dirección; revisiones de sprint para los dueños de producto; validaciones de su ámbito para TI, Prevención de Pérdidas y Contraloría; boletín mensual a jefaturas, semanal en marcha blanca; mensajes breves y demostraciones al personal de tienda; externos siempre a través de Ancoa. & Jefe de Proyecto \\ \addlinespace
\textbf{Capacitación y formación} & Personal de tienda, dueños de producto, TI, repositores externos y soporte de niveles 1 y 2 & Módulos de no más de 30 minutos, ambiente de entrenamiento, ayuda en pantalla e inducción permanente; formación presencial, sincrónica, autoformación y acompañamiento en puesto; material editable, en español y de propiedad de Ancoa; dos jornadas anuales de actualización en operación y mentoría técnica de al menos seis meses tras la producción de la Etapa 2. & Líder de Implantación y Gestión del Cambio \\ \addlinespace
\textbf{Participación y co-creación} & Dueños de producto, usuarios clave, tiendas del piloto y terceros & Revisiones de sprint, pruebas de aceptación, pilotos antes de cada ola y talleres de integración con terceros a través de Ancoa. & Líder Funcional \\ \addlinespace
\textbf{Reconocimiento y refuerzo positivo} & Tiendas, agentes de cambio, usuarios piloto y personal de venta & Destacar mejoras y metas de adopción alcanzadas; registrar los aportes en las notas de versión; entregar a la Gerencia Comercial y a Gestión de Personas los datos para evaluar ajustes a la regla de comisión, cuya decisión corresponde a Ancoa. & Líder de Implantación y Gestión del Cambio \\
\end{longtable}

\subsubsection{Tensiones entre interesados}\label{t9-4-tensiones}

El análisis de actores del Subdocumento 2 anticipa siete tensiones. Only Simple Solutions no las resolverá por su cuenta: las hará visibles, preparará alternativas con datos y llevará la decisión a la instancia que corresponda. La Tabla~\mbox{\ref{tab:t9-tensiones}} registra cada una con su tratamiento.

\begin{longtable}[]{@{}
  >{\raggedright\arraybackslash}p{(\columnwidth - 10\tabcolsep) * \real{0.0600}}
  >{\raggedright\arraybackslash}p{(\columnwidth - 10\tabcolsep) * \real{0.2400}}
  >{\raggedright\arraybackslash}p{(\columnwidth - 10\tabcolsep) * \real{0.1300}}
  >{\raggedright\arraybackslash}p{(\columnwidth - 10\tabcolsep) * \real{0.3500}}
  >{\raggedright\arraybackslash}p{(\columnwidth - 10\tabcolsep) * \real{0.2200}}@{}}
\caption{Tensiones entre interesados y su tratamiento.}\label{tab:t9-tensiones}\\
\toprule\noalign{}
\begin{minipage}[b]{\linewidth}\raggedright
\textbf{N.º}
\end{minipage} & \begin{minipage}[b]{\linewidth}\raggedright
\textbf{Tensión}
\end{minipage} & \begin{minipage}[b]{\linewidth}\raggedright
\textbf{Interesados}
\end{minipage} & \begin{minipage}[b]{\linewidth}\raggedright
\textbf{Tratamiento}
\end{minipage} & \begin{minipage}[b]{\linewidth}\raggedright
\textbf{Decide}
\end{minipage} \\
\midrule\noalign{}
\endfirsthead
\toprule\noalign{}
\begin{minipage}[b]{\linewidth}\raggedright
\textbf{N.º}
\end{minipage} & \begin{minipage}[b]{\linewidth}\raggedright
\textbf{Tensión}
\end{minipage} & \begin{minipage}[b]{\linewidth}\raggedright
\textbf{Interesados}
\end{minipage} & \begin{minipage}[b]{\linewidth}\raggedright
\textbf{Tratamiento}
\end{minipage} & \begin{minipage}[b]{\linewidth}\raggedright
\textbf{Decide}
\end{minipage} \\
\midrule\noalign{}
\endhead
\bottomrule\noalign{}
\multicolumn{5}{@{}l}{\rule{0pt}{3.2ex}\textbf{Fuente:} Elaboración propia.} \\
\endlastfoot
\textbf{1} & Rapidez del crédito frente a evidencia del consentimiento & I-09, I-03 & La evaluación y la entrega de la información precontractual se diseñan como un solo flujo. & Validación conjunta de I-09 e I-03 antes del paso a producción. \\ \addlinespace
\textbf{2} & Comisión del vendedor frente a los pasos de cumplimiento & I-14, I-03 & Se mide la duración del flujo de venta con crédito en las tiendas del piloto. & Gerencia Comercial y Gestión de Personas. \\ \addlinespace
\textbf{3} & Existencia publicada frente a existencia confiable & I-07, I-08 & La regla de publicación se define con los datos de exactitud de inventario de cada ola. & I-07 e I-08; sin acuerdo, Comité de Proyecto. \\ \addlinespace
\textbf{4} & Velocidad del cambio de precio frente al precio exhibido & I-06, I-13 & Los cambios de precio se programan en oleadas coordinadas con el reetiquetado. & Gerencia Comercial con las jefaturas; sin acuerdo, Comité de Proyecto. \\ \addlinespace
\textbf{5} & Despacho desde tienda frente a la venta en sala & I-07, I-13, I-14 & Se entregan datos de ambos canales sobre asignación de unidades y reconocimiento de la venta. & Gerencia Comercial. \\ \addlinespace
\textbf{6} & Uso comercial de datos financieros & I-06, I-03, I-09 & Ningún cruce sin finalidad definida, base de licitud y aprobación de Contraloría y del Negocio Financiero (Ley N.º 21.719, 2024). & Instancia competente, con las evaluaciones documentadas. \\ \addlinespace
\textbf{7} & Prioridad del inventario frente al plazo de 2029 & I-02, I-09 & La preparación de la migración de cartera se planifica desde la Etapa 1. & Comité Ejecutivo. \\
\end{longtable}

Cada conflicto se registra en Jira con su responsable de Ancoa y la fecha requerida de decisión; se analiza su impacto en alcance, plazo, costo y riesgo; se concilia en una reunión facilitada por el Scrum Master o por el Líder de Implantación; y su acuerdo se comunica y, si modifica una línea base, se tramita por control integrado de cambios. Si no se resuelve, sigue la ruta de escalamiento de la Tabla~\mbox{\ref{tab:t9-escalamiento}}.

\subsubsection{Gestión del cambio con ADKAR}\label{t9-4-adkar}

Para el grupo Operación de tienda, la participación se complementa con el modelo ADKAR (Prosci, s.f.). La Tabla~\mbox{\ref{tab:t9-adkar}} fija, para cada ola, cuándo se trabaja cada elemento y cómo se comprueba su logro.

\begin{longtable}[]{@{}
  >{\raggedright\arraybackslash}p{(\columnwidth - 8\tabcolsep) * \real{0.1600}}
  >{\raggedright\arraybackslash}p{(\columnwidth - 8\tabcolsep) * \real{0.2200}}
  >{\raggedright\arraybackslash}p{(\columnwidth - 8\tabcolsep) * \real{0.3400}}
  >{\raggedright\arraybackslash}p{(\columnwidth - 8\tabcolsep) * \real{0.2800}}@{}}
\caption{Secuencia ADKAR por ola de implantación.}\label{tab:t9-adkar}\\
\toprule\noalign{}
\begin{minipage}[b]{\linewidth}\raggedright
\textbf{Elemento}
\end{minipage} & \begin{minipage}[b]{\linewidth}\raggedright
\textbf{Inicio}
\end{minipage} & \begin{minipage}[b]{\linewidth}\raggedright
\textbf{Acciones}
\end{minipage} & \begin{minipage}[b]{\linewidth}\raggedright
\textbf{Evidencia de logro}
\end{minipage} \\
\midrule\noalign{}
\endfirsthead
\toprule\noalign{}
\begin{minipage}[b]{\linewidth}\raggedright
\textbf{Elemento}
\end{minipage} & \begin{minipage}[b]{\linewidth}\raggedright
\textbf{Inicio}
\end{minipage} & \begin{minipage}[b]{\linewidth}\raggedright
\textbf{Acciones}
\end{minipage} & \begin{minipage}[b]{\linewidth}\raggedright
\textbf{Evidencia de logro}
\end{minipage} \\
\midrule\noalign{}
\endhead
\bottomrule\noalign{}
\multicolumn{4}{@{}l}{\rule{0pt}{3.2ex}\textbf{Fuente:} Elaboración propia a partir de Prosci (s.f.).} \\
\endlastfoot
\textbf{Conciencia} & Ocho semanas antes de la ola & Mensajes de la gerencia responsable por medio de cada jefatura sobre el propósito del cambio. & Comprensión verificada en encuesta breve. \\ \addlinespace
\textbf{Deseo} & Ocho semanas antes de la ola & Agentes de cambio, reconocimiento y datos para la revisión de la comisión. & Participación en pilotos y demostraciones. \\ \addlinespace
\textbf{Conocimiento} & Cuatro semanas antes de la ola & Módulos breves y práctica en el ambiente de entrenamiento. & Aprobación de la capacitación. \\ \addlinespace
\textbf{Habilidad} & Cuatro semanas antes de la ola & Práctica guiada y acompañamiento en puesto. & Ejecución práctica observada. \\ \addlinespace
\textbf{Refuerzo} & Durante la marcha blanca y el mes siguiente & Ayuda en pantalla, sala de control y reconocimiento. & Uso sostenido y consultas resueltas. \\
\end{longtable}

El logro de un elemento no se da por alcanzado solo porque transcurra el tiempo. La secuencia comienza en las tiendas del piloto, y sus lecciones ajustan el plan antes de la primera ola.

\subsubsection{Verificación de participación antes de cada marcha blanca}\label{t9-4-verificacion}

Antes de cada marcha blanca y paso a producción, el Jefe de Proyecto y la contraparte técnica comprueban las cinco condiciones de participación siguientes.

\begin{enumerate}
\tightlist
\item
  Están representadas las áreas realmente afectadas, con representantes que tienen las atribuciones necesarias.
\item
  Los dueños de producto están designados y cumplen su dedicación, y no hay decisiones de Ancoa pendientes fuera de plazo.
\item
  La estrategia no vulnera la separación entre retail y emisor; las materias de crédito y datos tienen las validaciones de datos, seguridad y cumplimiento.
\item
  Las tiendas de la ola tienen su agente de cambio y su personal capacitado.
\item
  Los terceros afectados recibieron aviso con anticipación a través de Ancoa.
\end{enumerate}

Si una condición no se cumple, se registra como impedimento o riesgo y el Jefe de Proyecto acuerda una alternativa con la contraparte técnica. Una sesión con usuarios o una aprobación comercial no sustituye las validaciones de datos, seguridad y cumplimiento.

\subsubsection{Indicadores y acciones de mejora}\label{t9-4-indicadores}

Las metas de la Tabla~\mbox{\ref{tab:t9-ind-interesados}} son umbrales internos de gestión y no reducen obligaciones contractuales ni condiciones de aceptación. La tabla agrega, a los indicadores de la Tabla 6.5 del Subdocumento 6, el responsable de medirlos y la acción que se activa bajo la meta.

\begin{longtable}[]{@{}
  >{\raggedright\arraybackslash}p{(\columnwidth - 8\tabcolsep) * \real{0.2600}}
  >{\raggedright\arraybackslash}p{(\columnwidth - 8\tabcolsep) * \real{0.1800}}
  >{\raggedright\arraybackslash}p{(\columnwidth - 8\tabcolsep) * \real{0.2200}}
  >{\raggedright\arraybackslash}p{(\columnwidth - 8\tabcolsep) * \real{0.3400}}@{}}
\caption{Indicadores de involucramiento, responsable y acción.}\label{tab:t9-ind-interesados}\\
\toprule\noalign{}
\begin{minipage}[b]{\linewidth}\raggedright
\textbf{Indicador}
\end{minipage} & \begin{minipage}[b]{\linewidth}\raggedright
\textbf{Meta}
\end{minipage} & \begin{minipage}[b]{\linewidth}\raggedright
\textbf{Responsable}
\end{minipage} & \begin{minipage}[b]{\linewidth}\raggedright
\textbf{Acción bajo la meta}
\end{minipage} \\
\midrule\noalign{}
\endfirsthead
\toprule\noalign{}
\begin{minipage}[b]{\linewidth}\raggedright
\textbf{Indicador}
\end{minipage} & \begin{minipage}[b]{\linewidth}\raggedright
\textbf{Meta}
\end{minipage} & \begin{minipage}[b]{\linewidth}\raggedright
\textbf{Responsable}
\end{minipage} & \begin{minipage}[b]{\linewidth}\raggedright
\textbf{Acción bajo la meta}
\end{minipage} \\
\midrule\noalign{}
\endhead
\bottomrule\noalign{}
\multicolumn{4}{@{}l}{\rule{0pt}{3.2ex}\textbf{Fuente:} Elaboración propia a partir de PMI (2017).} \\
\endlastfoot
\textbf{Participación efectiva} & ≥ 85 \% & Jefe de Proyecto & Reprogramar sesiones en horarios de menor afluencia o cambiar el canal. \\ \addlinespace
\textbf{Cumplimiento de compromisos} & ≥ 90 \% & Jefe de Proyecto & Revisar compromisos vencidos en el Comité de Proyecto. \\ \addlinespace
\textbf{Decisiones de Ancoa a tiempo} & ≥ 90 \% & Jefe de Proyecto & Solicitar a la gerencia el cambio o la incorporación de representantes. \\ \addlinespace
\textbf{Satisfacción de los interesados} & ≥ 4,0 en escala de 1 a 5 & Líder de Implantación & Ajustar formato y frecuencia de las revisiones y pilotos. \\ \addlinespace
\textbf{Resolución de conflictos} & ≤ 10 días hábiles & Jefe de Proyecto & Escalar según la ruta de escalamiento. \\ \addlinespace
\textbf{Cobertura de agentes de cambio} & 22 de 22 antes de H6 & Líder de Implantación & Solicitar la designación a las jefaturas pendientes. \\ \addlinespace
\textbf{Personal capacitado por ola} & ≥ 95 \% antes de cada ola & Líder de Implantación & Activar un plan de recuperación de la ola, sin dispensar capacitación ni certificación. \\ \addlinespace
\textbf{Interesados internos críticos en nivel deseado} & 100 \% antes de cada paso a producción & Jefe de Proyecto & Reuniones extraordinarias con el área y, si persiste, escalamiento al Comité Ejecutivo. \\
\end{longtable}

Se activa una acción de mejora cuando un indicador queda bajo su meta en dos mediciones seguidas o cuando un interesado de gestión cercana retrocede un nivel en la matriz. Si un interesado crítico interno permanece bajo el nivel requerido durante dos intervalos seguidos, el asunto se escala al Comité Ejecutivo. Si no existe población aplicable en un período, se registra esa situación en lugar de un porcentaje sin denominador.

\subsection{T-9.5 Gestión de comunicaciones}\label{t9-5}

\subsubsection{Sistema de información para la dirección del proyecto}\label{t9-5-pmis}

Las comunicaciones se apoyan en una fuente única de datos de gestión. La Tabla~\mbox{\ref{tab:t9-pmis}} define la función de cada herramienta, lo que registra y quién accede.

\begin{longtable}[]{@{}
  >{\raggedright\arraybackslash}p{(\columnwidth - 8\tabcolsep) * \real{0.1800}}
  >{\raggedright\arraybackslash}p{(\columnwidth - 8\tabcolsep) * \real{0.2200}}
  >{\raggedright\arraybackslash}p{(\columnwidth - 8\tabcolsep) * \real{0.3400}}
  >{\raggedright\arraybackslash}p{(\columnwidth - 8\tabcolsep) * \real{0.2600}}@{}}
\caption{Herramientas del PMIS y del canal de coordinación.}\label{tab:t9-pmis}\\
\toprule\noalign{}
\begin{minipage}[b]{\linewidth}\raggedright
\textbf{Herramienta}
\end{minipage} & \begin{minipage}[b]{\linewidth}\raggedright
\textbf{Función}
\end{minipage} & \begin{minipage}[b]{\linewidth}\raggedright
\textbf{Qué registra}
\end{minipage} & \begin{minipage}[b]{\linewidth}\raggedright
\textbf{Acceso}
\end{minipage} \\
\midrule\noalign{}
\endfirsthead
\toprule\noalign{}
\begin{minipage}[b]{\linewidth}\raggedright
\textbf{Herramienta}
\end{minipage} & \begin{minipage}[b]{\linewidth}\raggedright
\textbf{Función}
\end{minipage} & \begin{minipage}[b]{\linewidth}\raggedright
\textbf{Qué registra}
\end{minipage} & \begin{minipage}[b]{\linewidth}\raggedright
\textbf{Acceso}
\end{minipage} \\
\midrule\noalign{}
\endhead
\bottomrule\noalign{}
\multicolumn{4}{@{}l}{\rule{0pt}{3.2ex}\textbf{Fuente:} Elaboración propia.} \\
\endlastfoot
\textbf{Jira} & Registro operativo del PMIS & Tareas, responsables, plazos, incidencias, compromisos, dependencias y decisiones; tablero de estado del proyecto. & Equipos y representantes autorizados de Ancoa, según su rol. \\ \addlinespace
\textbf{Microsoft SharePoint} & Espacio colaborativo y repositorio documental del PMIS & Entregables, actas, registro de riesgos, registro de cambios, documentación vigente y evidencias, con responsable, fecha y versión. & Permisos por función y clasificación; la información del negocio financiero, restringida a personas autorizadas. \\ \addlinespace
\textbf{Slack} & Coordinación cotidiana & Consultas e impedimentos del día a día. No constituye registro formal. & Integrantes de los equipos del proyecto. \\
\end{longtable}

Todo acuerdo de Slack que afecte decisiones, responsabilidades, plazos o entregables se formaliza en Jira o SharePoint. Con 26 interesados existen, como máximo teórico, $n(n-1)/2 = 26 \times 25 / 2 = 325$ pares de comunicación; esta cifra justifica canales y niveles definidos en vez de una comunicación indiscriminada, y no representa conversaciones efectivas.

\subsubsection{Matriz de comunicaciones}\label{t9-5-matriz}

La Tabla~\mbox{\ref{tab:t9-matriz-com}} es la matriz completa de comunicaciones. Cada fila indica quién emite, quién recibe, el método de comunicación (interactiva, dirigida o disponible) y la oportunidad. Las cuatro últimas filas son comunicaciones extraordinarias, que se emiten sin esperar la siguiente reunión programada.

\begin{longtable}[]{@{}
  >{\raggedright\arraybackslash}p{(\columnwidth - 10\tabcolsep) * \real{0.2400}}
  >{\raggedright\arraybackslash}p{(\columnwidth - 10\tabcolsep) * \real{0.2000}}
  >{\raggedright\arraybackslash}p{(\columnwidth - 10\tabcolsep) * \real{0.2200}}
  >{\raggedright\arraybackslash}p{(\columnwidth - 10\tabcolsep) * \real{0.1400}}
  >{\raggedright\arraybackslash}p{(\columnwidth - 10\tabcolsep) * \real{0.2000}}@{}}
\caption{Matriz de comunicaciones del proyecto.}\label{tab:t9-matriz-com}\\
\toprule\noalign{}
\begin{minipage}[b]{\linewidth}\raggedright
\textbf{Comunicación}
\end{minipage} & \begin{minipage}[b]{\linewidth}\raggedright
\textbf{Emisor}
\end{minipage} & \begin{minipage}[b]{\linewidth}\raggedright
\textbf{Destinatario}
\end{minipage} & \begin{minipage}[b]{\linewidth}\raggedright
\textbf{Método}
\end{minipage} & \begin{minipage}[b]{\linewidth}\raggedright
\textbf{Oportunidad}
\end{minipage} \\
\midrule\noalign{}
\endfirsthead
\toprule\noalign{}
\begin{minipage}[b]{\linewidth}\raggedright
\textbf{Comunicación}
\end{minipage} & \begin{minipage}[b]{\linewidth}\raggedright
\textbf{Emisor}
\end{minipage} & \begin{minipage}[b]{\linewidth}\raggedright
\textbf{Destinatario}
\end{minipage} & \begin{minipage}[b]{\linewidth}\raggedright
\textbf{Método}
\end{minipage} & \begin{minipage}[b]{\linewidth}\raggedright
\textbf{Oportunidad}
\end{minipage} \\
\midrule\noalign{}
\endhead
\bottomrule\noalign{}
\multicolumn{5}{@{}l}{\rule{0pt}{3.2ex}\textbf{Fuente:} Elaboración propia a partir de PMI (2017).} \\
\endlastfoot
\textbf{Sesión del Comité Ejecutivo} & Jefe de Proyecto & Patrocinador, gerencia de OSS y administrador del contrato & Interactiva & Mensual \\ \addlinespace
\textbf{Sesión del Comité de Proyecto} & Jefe de Proyecto & Contraparte técnica y dueños de producto según agenda & Interactiva & Quincenal \\ \addlinespace
\textbf{Sesión del Comité de Arquitectura} & Arquitecto de Solución & Encargado de Seguridad y referentes técnicos de Ancoa & Interactiva & Mensual \\ \addlinespace
\textbf{Sesión del Comité de Operación} & Líder de Operación / SRE & Mesa de servicio y contraparte técnica & Interactiva & Mensual desde el mes 13 \\ \addlinespace
\textbf{Actas de los cuatro comités} & Jefe de Proyecto, con el responsable designado & Participantes autorizados & Dirigida & Dentro de dos días hábiles \\ \addlinespace
\textbf{Reunión de seguimiento} & Jefe de Proyecto & Equipos de ambas partes & Interactiva & Semanal \\ \addlinespace
\textbf{Informe de desempeño} & Jefe de Proyecto & Contrapartes y Comité Ejecutivo & Dirigida & Mensual \\ \addlinespace
\textbf{Tablero de estado del proyecto} & Jefe de Proyecto & Representantes autorizados de Ancoa & Disponible & Datos operativos diarios; consolidados con fecha de corte \\ \addlinespace
\textbf{Daily Scrum} & Developers & Equipo Scrum & Interactiva & Diaria, 15 minutos \\ \addlinespace
\textbf{Revisión del sprint} & Equipo de desarrollo & Dueño de producto y usuarios clave & Interactiva & Cada dos semanas \\ \addlinespace
\textbf{Coordinación técnica de dependencias} & Arquitecto de Solución y responsables de desarrollo & Equipos y TI de Ancoa & Interactiva & Según el tablero de dependencias \\ \addlinespace
\textbf{Contrato de interfaz} & Responsable técnico de la integración & Equipos y terceros, a través de Ancoa & Disponible & Antes de construir la integración \\ \addlinespace
\textbf{Decisión de arquitectura aprobada} & Arquitecto de Solución & Equipos y TI de Ancoa & Disponible & Al aprobarse en el comité \\ \addlinespace
\textbf{Boletín a tiendas} & Líder de Implantación; anuncia la gerencia responsable & Jefaturas y agentes de cambio & Dirigida & Mensual; semanal en marcha blanca \\ \addlinespace
\textbf{Mensajes breves y demostraciones} & Gerencia responsable, por medio de cada jefatura & Personal de tienda & Dirigida & Por ola \\ \addlinespace
\textbf{Cápsulas de capacitación y ayuda} & Líder de Implantación & Personal de tienda & Disponible & Permanente \\ \addlinespace
\textbf{Sala de control de marcha blanca} & Líder de Operación y Líder de Implantación & Tiendas de la ola y contraparte técnica & Interactiva & Diaria durante la marcha blanca \\ \addlinespace
\textbf{Nota técnica y programa de trabajos} & Arquitecto de Solución y Líder de Implantación & Terceros afectados, a través de Ancoa & Dirigida & Antes del cambio o del trabajo físico \\ \addlinespace
\textbf{Resumen para Propiedad} & Jefe de Proyecto, por medio del Directorio & I-04 e I-05 & Dirigida & Por hito \\ \addlinespace
\textbf{Evidencia regulatoria y de frontera} & Responsables de datos y cumplimiento, con seguridad & Filial o gerencia de Ancoa que responde & Dirigida & Según el procedimiento aplicable \\ \addlinespace
\textbf{Comunicación de migración de cartera} & Líder de Datos y Líder Funcional & Filial emisora y responsables de atención; titulares por ola & Dirigida & Por ola de migración \\ \addlinespace
\textbf{Desviación de hito superior al 10 \%} & Jefe de Proyecto & Administrador del contrato y contraparte técnica & Dirigida & Aviso y plan dentro de cinco días hábiles \\ \addlinespace
\textbf{Incidente crítico de seguridad} & Encargado de Seguridad & Contrapartes de seguridad y técnica & Dirigida & Dentro de dos horas desde la detección \\ \addlinespace
\textbf{Brecha de seguridad o de datos personales} & Encargado de Seguridad & Contrapartes autorizadas & Dirigida & Informe preliminar dentro de 24 horas \\ \addlinespace
\textbf{Causa raíz de la brecha} & Encargado de Seguridad & Contrapartes autorizadas & Dirigida & Dentro de cinco días hábiles \\
\end{longtable}

Ninguna comunicación de Only Simple Solutions sustituye la representación institucional de Ancoa o de su filial ante una autoridad. Para las comunicaciones extraordinarias se mantienen contactos titulares y suplentes y se registra la fecha de detección, de emisión y de recepción de cada aviso. Los avisos de incidente crítico y de brecha son controles distintos y se cumplen conjuntamente cuando ambos aplican.

\subsubsection{Contenidos mínimos}\label{t9-5-contenidos}

Para que cada comunicación sea verificable, la Tabla~\mbox{\ref{tab:t9-contenidos}} fija su contenido mínimo.

\begin{longtable}[]{@{}
  >{\raggedright\arraybackslash}p{(\columnwidth - 4\tabcolsep) * \real{0.2600}}
  >{\raggedright\arraybackslash}p{(\columnwidth - 4\tabcolsep) * \real{0.7400}}@{}}
\caption{Contenido mínimo de las comunicaciones formales.}\label{tab:t9-contenidos}\\
\toprule\noalign{}
\begin{minipage}[b]{\linewidth}\raggedright
\textbf{Comunicación}
\end{minipage} & \begin{minipage}[b]{\linewidth}\raggedright
\textbf{Contenido mínimo}
\end{minipage} \\
\midrule\noalign{}
\endfirsthead
\toprule\noalign{}
\begin{minipage}[b]{\linewidth}\raggedright
\textbf{Comunicación}
\end{minipage} & \begin{minipage}[b]{\linewidth}\raggedright
\textbf{Contenido mínimo}
\end{minipage} \\
\midrule\noalign{}
\endhead
\bottomrule\noalign{}
\multicolumn{2}{@{}l}{\rule{0pt}{3.2ex}\textbf{Fuente:} Elaboración propia.} \\
\endlastfoot
\textbf{Informe mensual de desempeño} & Estado del cronograma y de los hitos; avance físico y financiero; valor ganado con SPI y CPI y su fecha de corte; entregables del período; desviaciones y su explicación; riesgos; incidencias; decisiones pendientes; compromisos del período siguiente. \\ \addlinespace
\textbf{Acta de comité} & Fecha, asistentes y quórum de representantes autorizados; decisiones adoptadas; acuerdos con responsable y plazo; asuntos que siguen abiertos. \\ \addlinespace
\textbf{Aviso de desviación de hito} & Fecha de detección; hito afectado; magnitud de la desviación; causas; consecuencias sobre las dependencias; plan de recuperación con medidas, responsables, plazos y forma de verificar la recuperación. \\ \addlinespace
\textbf{Aviso inicial de incidente crítico} & Fecha y hora de detección; sistemas o servicios potencialmente afectados; severidad conocida; medidas de contención iniciadas; responsable de coordinar la respuesta. Se actualiza con cada antecedente relevante; la convocatoria de la mesa de análisis no suspende el plazo. \\ \addlinespace
\textbf{Informe de brecha} & Informe preliminar dentro de 24 horas y análisis de causa raíz dentro de cinco días hábiles, con la información del negocio financiero tratada solo con las contrapartes autorizadas de la filial. \\ \addlinespace
\textbf{Mensaje a tiendas} & Qué cambia, cuándo cambia y a quién preguntar. \\
\end{longtable}

\subsubsection{Monitoreo y control de las comunicaciones}\label{t9-5-monitoreo}

El monitoreo verifica que la información llegue, se entienda y produzca las decisiones esperadas. Se ejecuta como un ciclo de seis fases por intervalo de planificación, más corto durante la marcha blanca. La Tabla~\mbox{\ref{tab:t9-ciclo}} detalla cada fase.

\begin{longtable}[]{@{}
  >{\raggedright\arraybackslash}p{(\columnwidth - 8\tabcolsep) * \real{0.1600}}
  >{\raggedright\arraybackslash}p{(\columnwidth - 8\tabcolsep) * \real{0.4400}}
  >{\raggedright\arraybackslash}p{(\columnwidth - 8\tabcolsep) * \real{0.1800}}
  >{\raggedright\arraybackslash}p{(\columnwidth - 8\tabcolsep) * \real{0.2200}}@{}}
\caption{Ciclo de monitoreo de las comunicaciones.}\label{tab:t9-ciclo}\\
\toprule\noalign{}
\begin{minipage}[b]{\linewidth}\raggedright
\textbf{Fase}
\end{minipage} & \begin{minipage}[b]{\linewidth}\raggedright
\textbf{Actividad}
\end{minipage} & \begin{minipage}[b]{\linewidth}\raggedright
\textbf{Responsable}
\end{minipage} & \begin{minipage}[b]{\linewidth}\raggedright
\textbf{Salida}
\end{minipage} \\
\midrule\noalign{}
\endfirsthead
\toprule\noalign{}
\begin{minipage}[b]{\linewidth}\raggedright
\textbf{Fase}
\end{minipage} & \begin{minipage}[b]{\linewidth}\raggedright
\textbf{Actividad}
\end{minipage} & \begin{minipage}[b]{\linewidth}\raggedright
\textbf{Responsable}
\end{minipage} & \begin{minipage}[b]{\linewidth}\raggedright
\textbf{Salida}
\end{minipage} \\
\midrule\noalign{}
\endhead
\bottomrule\noalign{}
\multicolumn{4}{@{}l}{\rule{0pt}{3.2ex}\textbf{Fuente:} Elaboración propia a partir de PMI (2017).} \\
\endlastfoot
\textbf{1. Recolectar} & Reunir asistencia, actas y su fecha de publicación, acuses de lectura, uso del tablero y del repositorio, resultados de la encuesta de eficacia y fechas que acreditan cada plazo. & Responsable de cada frente & Datos con fuente y evidencia en el PMIS. \\ \addlinespace
\textbf{2. Medir} & Comprobar por evento los plazos obligatorios y, por separado, las metas internas de eficacia. & Jefe de Proyecto & Indicadores del período. \\ \addlinespace
\textbf{3. Analizar} & Identificar si el problema está en el contenido, el canal, la frecuencia o el destinatario; en un plazo incumplido, revisar detección, responsabilidad y escalamiento. & Responsable del indicador & Causa identificada. \\ \addlinespace
\textbf{4. Decidir} & Presentar desviaciones y acciones correctivas al Comité de Proyecto, que las aprueba con responsable y plazo. & Jefe de Proyecto & Acciones aprobadas. \\ \addlinespace
\textbf{5. Ajustar} & Aplicar las acciones y actualizar la matriz, los canales o la frecuencia; informar a los afectados. & Jefe de Proyecto & Nueva versión del plan, si el ajuste es relevante. \\ \addlinespace
\textbf{6. Registrar} & Anotar cada hallazgo con fecha, canal, observación y acción correctiva. & Jefe de Proyecto & Bitácora y registro de lecciones aprendidas. \\
\end{longtable}

La Tabla~\mbox{\ref{tab:t9-ind-com}} reúne los indicadores del monitoreo. Los plazos obligatorios no admiten tolerancia; las demás metas son internas.

\begin{longtable}[]{@{}
  >{\raggedright\arraybackslash}p{(\columnwidth - 6\tabcolsep) * \real{0.3600}}
  >{\raggedright\arraybackslash}p{(\columnwidth - 6\tabcolsep) * \real{0.2000}}
  >{\raggedright\arraybackslash}p{(\columnwidth - 6\tabcolsep) * \real{0.4400}}@{}}
\caption{Indicadores del monitoreo de las comunicaciones.}\label{tab:t9-ind-com}\\
\toprule\noalign{}
\begin{minipage}[b]{\linewidth}\raggedright
\textbf{Indicador}
\end{minipage} & \begin{minipage}[b]{\linewidth}\raggedright
\textbf{Meta}
\end{minipage} & \begin{minipage}[b]{\linewidth}\raggedright
\textbf{Medición}
\end{minipage} \\
\midrule\noalign{}
\endfirsthead
\toprule\noalign{}
\begin{minipage}[b]{\linewidth}\raggedright
\textbf{Indicador}
\end{minipage} & \begin{minipage}[b]{\linewidth}\raggedright
\textbf{Meta}
\end{minipage} & \begin{minipage}[b]{\linewidth}\raggedright
\textbf{Medición}
\end{minipage} \\
\midrule\noalign{}
\endhead
\bottomrule\noalign{}
\multicolumn{3}{@{}l}{\rule{0pt}{3.2ex}\textbf{Fuente:} Elaboración propia.} \\
\endlastfoot
\textbf{Cumplimiento de plazos obligatorios} & 100 \% por evento & Actas en dos días hábiles, aviso de desviación en cinco días hábiles, incidente crítico en dos horas, brecha en 24 horas y causa raíz en cinco días hábiles. \\ \addlinespace
\textbf{Lectura del boletín por agentes de cambio} & ≥ 90 \% & Acuses de lectura por envío. \\ \addlinespace
\textbf{Participación interna} & ≥ 85 \% & Listas de asistencia de comités, revisiones y talleres. \\ \addlinespace
\textbf{Claridad de la comunicación} & ≥ 4 en escala de 1 a 5 & Encuesta breve de eficacia por intervalo. \\ \addlinespace
\textbf{Trazabilidad de acuerdos} & 100 \% & Acuerdos de acta con responsable y plazo vinculados en Jira. \\
\end{longtable}

El Scrum Master revisa a diario la comunicación de su equipo; el Jefe de Proyecto revisa cada semana actas, acuerdos y bitácora; los indicadores se incluyen cada mes en el informe de desempeño; y al cierre de cada intervalo se aplica la encuesta de eficacia. Durante la marcha blanca, la sala de control revisa a diario la comunicación con las tiendas. Si una meta interna permanece incumplida en dos revisiones consecutivas, el Jefe de Proyecto la escala al Comité Ejecutivo; un plazo obligatorio incumplido se atiende desde su detección. Al cierre de cada etapa se entrega una evaluación de comunicaciones con medidas aplicables a la operación.

\subsection{T-9.6 Gestión de adquisiciones y dependencias externas}\label{t9-6}

\subsubsection{Registro de dependencias externas}\label{t9-6-registro}

Cada dependencia externa se registra en Jira, vinculada con la actividad o el entregable que condiciona, antes de comprometer el trabajo que la requiere. La Tabla~\mbox{\ref{tab:t9-dep-campos}} fija los campos del registro.

\begin{longtable}[]{@{}
  >{\raggedright\arraybackslash}p{(\columnwidth - 4\tabcolsep) * \real{0.3000}}
  >{\raggedright\arraybackslash}p{(\columnwidth - 4\tabcolsep) * \real{0.7000}}@{}}
\caption{Campos del registro de dependencias externas.}\label{tab:t9-dep-campos}\\
\toprule\noalign{}
\begin{minipage}[b]{\linewidth}\raggedright
\textbf{Campo}
\end{minipage} & \begin{minipage}[b]{\linewidth}\raggedright
\textbf{Contenido}
\end{minipage} \\
\midrule\noalign{}
\endfirsthead
\toprule\noalign{}
\begin{minipage}[b]{\linewidth}\raggedright
\textbf{Campo}
\end{minipage} & \begin{minipage}[b]{\linewidth}\raggedright
\textbf{Contenido}
\end{minipage} \\
\midrule\noalign{}
\endhead
\bottomrule\noalign{}
\multicolumn{2}{@{}l}{\rule{0pt}{3.2ex}\textbf{Fuente:} Elaboración propia.} \\
\endlastfoot
\textbf{Identificador y tipo} & Código único y tipo de dependencia según la Tabla~\mbox{\ref{tab:t9-dep-tipos}}. \\ \addlinespace
\textbf{Resultado esperado} & Componente, servicio, acceso, habilitación o trabajo que debe estar disponible. \\ \addlinespace
\textbf{Especificación técnica} & Cantidades, capacidades, características, interfaces, compatibilidad, instalación, soporte, continuidad y administración. \\ \addlinespace
\textbf{Responsable y contraparte} & Responsable dentro de Only Simple Solutions y contraparte de Ancoa o del proveedor. \\ \addlinespace
\textbf{Trabajo afectado y fecha requerida} & Paquete de trabajo o entregable que condiciona y fecha en que debe estar disponible. \\ \addlinespace
\textbf{Condiciones de entrada} & Condiciones que deben cumplirse antes de utilizar el resultado. \\ \addlinespace
\textbf{Criterio de recepción y evidencia} & Criterio definido antes de solicitar, adquirir o contratar, y evidencia que lo comprueba. \\ \addlinespace
\textbf{Riesgos y contingencia} & Riesgos asociados y, cuando la criticidad lo justifica, alternativa de contingencia. \\ \addlinespace
\textbf{Estado y seguimiento} & Estado vigente, impedimentos, compromisos abiertos y próxima acción con responsable y fecha. \\
\end{longtable}

El estado de cada dependencia recorre la secuencia siguiente.

\begin{enumerate}
\tightlist
\item
  Identificada, especificada y solicitada.
\item
  En provisión o ejecución.
\item
  Recibida y en verificación.
\item
  Aceptada y cerrada.
\item
  Afectada por una desviación, estado excepcional que registra la condición no cumplida, lo expuesto, la acción requerida, su responsable y su fecha.
\end{enumerate}

Ninguna actividad crítica se mantiene con un estado de “pendiente del proveedor”. Toda dependencia abierta tiene responsable, fecha requerida, impacto identificado y acción de seguimiento. Las dependencias que puedan afectar actividades próximas se revisan en la reunión semanal de seguimiento.

\subsubsection{Catálogo de dependencias por tipo}\label{t9-6-tipos}

La Tabla~\mbox{\ref{tab:t9-dep-tipos}} identifica, para cada tipo de dependencia, quién provee o ejecuta, la responsabilidad que conserva Only Simple Solutions y la instancia de escalamiento. Los criterios de recepción por tipo se resumen en la Tabla 6.8 del Subdocumento 6, y marca, modelo y soporte de los componentes se declaran en el Formulario T-11.

\begin{longtable}[]{@{}
  >{\raggedright\arraybackslash}p{(\columnwidth - 8\tabcolsep) * \real{0.2400}}
  >{\raggedright\arraybackslash}p{(\columnwidth - 8\tabcolsep) * \real{0.2400}}
  >{\raggedright\arraybackslash}p{(\columnwidth - 8\tabcolsep) * \real{0.3200}}
  >{\raggedright\arraybackslash}p{(\columnwidth - 8\tabcolsep) * \real{0.2000}}@{}}
\caption{Catálogo de dependencias externas por tipo.}\label{tab:t9-dep-tipos}\\
\toprule\noalign{}
\begin{minipage}[b]{\linewidth}\raggedright
\textbf{Tipo}
\end{minipage} & \begin{minipage}[b]{\linewidth}\raggedright
\textbf{Provee o ejecuta}
\end{minipage} & \begin{minipage}[b]{\linewidth}\raggedright
\textbf{Responsabilidad de OSS}
\end{minipage} & \begin{minipage}[b]{\linewidth}\raggedright
\textbf{Escalamiento}
\end{minipage} \\
\midrule\noalign{}
\endfirsthead
\toprule\noalign{}
\begin{minipage}[b]{\linewidth}\raggedright
\textbf{Tipo}
\end{minipage} & \begin{minipage}[b]{\linewidth}\raggedright
\textbf{Provee o ejecuta}
\end{minipage} & \begin{minipage}[b]{\linewidth}\raggedright
\textbf{Responsabilidad de OSS}
\end{minipage} & \begin{minipage}[b]{\linewidth}\raggedright
\textbf{Escalamiento}
\end{minipage} \\
\midrule\noalign{}
\endhead
\bottomrule\noalign{}
\multicolumn{4}{@{}l}{\rule{0pt}{3.2ex}\textbf{Fuente:} Elaboración propia.} \\
\endlastfoot
\textbf{Servicios Microsoft Azure} & Only Simple Solutions, con Microsoft o un integrador autorizado & Selección, configuración versionada, observabilidad, respaldo, reversibilidad e incorporación a la solución. & Comité de Arquitectura \\ \addlinespace
\textbf{Seguridad y conectividad (Fortinet o equivalente justificado)} & Only Simple Solutions; instalación física por Ancoa o su integrador cuando corresponda & Especificación, coordinación y verificación de segmentación, accesos y trazabilidad. & Comité de Arquitectura \\ \addlinespace
\textbf{Componentes on-premise de la solución y su licenciamiento} & Only Simple Solutions (Dell Technologies u opción técnicamente equivalente) & Provisión, configuración, recepción y pruebas en el entorno objetivo. & Comité de Arquitectura \\ \addlinespace
\textbf{Hardware de terreno y dispositivos de tiendas y centros de distribución} & Ancoa & Especificación, dimensionamiento, coordinación y verificación de su incorporación. & Comité de Proyecto \\ \addlinespace
\textbf{Obra civil de separación de instalaciones} & Ancoa & Condiciones técnicas del sitio y verificación de que permite iniciar la actividad siguiente. & Comité de Proyecto \\ \addlinespace
\textbf{Conectividad, seguridad y canalizaciones del centro local} & Only Simple Solutions & Provisión y habilitación. & Comité de Proyecto \\ \addlinespace
\textbf{Obras eléctricas, canalizaciones y cableado de tiendas y centros de distribución} & Según las exclusiones contractuales y el reparto de suministro por sitio & Condiciones técnicas por sitio y verificación de preparación. & Comité de Proyecto \\ \addlinespace
\textbf{Proveedores tecnológicos existentes de Ancoa} & Seis proveedores de nueve plataformas; gestión contractual de Ancoa & Dependencia técnica: especificar, preparar, seguir y advertir impedimentos de las catorce interfaces. & Comité de Proyecto; Comité Ejecutivo si excede atribuciones \\ \addlinespace
\textbf{Terceros operacionales} & Vendedores de marketplace, última milla, repositores externos y administradores de centros comerciales & Interfaces y requisitos técnicos, operacionales y de acceso, sin dirigir su operación ni su personal. & Comité de Proyecto \\ \addlinespace
\textbf{Subcontratistas y aliados} & Only Simple Solutions, con autorización previa y escrita de Ancoa y hasta el 40 \% del valor del contrato & Alcance, productos, fechas y controles iguales a los del equipo propio, sin transferir la responsabilidad. & Comité de Proyecto \\
\end{longtable}

Cuando la adquisición es de Ancoa, Only Simple Solutions entrega con anticipación las especificaciones necesarias para incorporarla. Un equipo no se considera habilitado solo por haber sido recibido, ni una obra disponible solo porque su ejecutor la informe terminada. La recepción técnica de una dependencia no equivale a la aceptación del entregable contractual en que participa.

\subsubsection{Responsabilidades en adquisiciones}\label{t9-6-resp}

La Tabla~\mbox{\ref{tab:t9-dep-roles}} asigna a cada rol de Only Simple Solutions su responsabilidad en la gestión de adquisiciones y dependencias externas.

\begin{longtable}[]{@{}
  >{\raggedright\arraybackslash}p{(\columnwidth - 4\tabcolsep) * \real{0.3000}}
  >{\raggedright\arraybackslash}p{(\columnwidth - 4\tabcolsep) * \real{0.7000}}@{}}
\caption{Responsabilidades en adquisiciones y dependencias.}\label{tab:t9-dep-roles}\\
\toprule\noalign{}
\begin{minipage}[b]{\linewidth}\raggedright
\textbf{Rol}
\end{minipage} & \begin{minipage}[b]{\linewidth}\raggedright
\textbf{Responsabilidad}
\end{minipage} \\
\midrule\noalign{}
\endfirsthead
\toprule\noalign{}
\begin{minipage}[b]{\linewidth}\raggedright
\textbf{Rol}
\end{minipage} & \begin{minipage}[b]{\linewidth}\raggedright
\textbf{Responsabilidad}
\end{minipage} \\
\midrule\noalign{}
\endhead
\bottomrule\noalign{}
\multicolumn{2}{@{}l}{\rule{0pt}{3.2ex}\textbf{Fuente:} Elaboración propia.} \\
\endlastfoot
\textbf{Jefe de Proyecto} & Consolida las dependencias en la planificación general, verifica su coherencia con el cronograma y coordina las que afectan hitos o requieren escalamiento. \\ \addlinespace
\textbf{Arquitecto de Solución} & Evaluación arquitectónica de servicios de nube, infraestructura y componentes: compatibilidad, interoperabilidad, escalabilidad, mantenibilidad y evolución. \\ \addlinespace
\textbf{Líder de Integración} & Dependencias de interfaz con las plataformas y los proveedores existentes. \\ \addlinespace
\textbf{Líder Funcional} & Reglas de negocio y condiciones funcionales asociadas. \\ \addlinespace
\textbf{Encargado de Seguridad de la Información} & Controles de acceso, segregación, protección de datos, trazabilidad y gestión de vulnerabilidades. \\ \addlinespace
\textbf{Líder de Datos} & Migración, almacenamiento, intercambio, transformación y conciliación, con atención a la separación entre los ámbitos comercial y financiero. \\ \addlinespace
\textbf{Líder de Calidad y Pruebas} & Criterios y evidencias suficientes antes de incorporar un resultado externo a un entregable. \\ \addlinespace
\textbf{Líder de Operación / SRE} & Condiciones de soporte, monitoreo, continuidad, recuperación y mantenimiento. \\ \addlinespace
\textbf{Líder de Implantación y Gestión del Cambio} & Preparación de sitios, disponibilidad de equipamiento y habilitaciones para los despliegues. \\
\end{longtable}

\subsubsection{Cierre y salida de terceros}\label{t9-6-cierre}

La participación de un proveedor, subcontratista o aliado se cierra solo cuando se cumplen las condiciones siguientes.

\begin{enumerate}
\tightlist
\item
  Se entregaron la documentación, las configuraciones, el código o los scripts que forman parte de la solución.
\item
  Los elementos correspondientes se incorporaron al inventario.
\item
  Las observaciones están resueltas y las incidencias abiertas, cerradas.
\item
  Se completó la transferencia de conocimiento.
\item
  Se revocaron los accesos y credenciales temporales, y se verificó la devolución o eliminación controlada de la información utilizada.
\item
  En servicios de nube y administrados, se conservan los datos, configuraciones, respaldos, definiciones de infraestructura y antecedentes necesarios para la continuidad o una transición futura.
\end{enumerate}

\subsection{T-9.7 Gobierno por etapas y transición a operación}\label{t9-7}

El énfasis de las instancias de gobierno evoluciona con el proyecto. La Tabla~\mbox{\ref{tab:t9-etapas}} indica, por período contractual, la fase en curso, el énfasis del gobierno y las instancias o roles que se incorporan.

\begin{longtable}[]{@{}
  >{\raggedright\arraybackslash}p{(\columnwidth - 8\tabcolsep) * \real{0.1400}}
  >{\raggedright\arraybackslash}p{(\columnwidth - 8\tabcolsep) * \real{0.2600}}
  >{\raggedright\arraybackslash}p{(\columnwidth - 8\tabcolsep) * \real{0.3200}}
  >{\raggedright\arraybackslash}p{(\columnwidth - 8\tabcolsep) * \real{0.2800}}@{}}
\caption{Gobierno por período contractual.}\label{tab:t9-etapas}\\
\toprule\noalign{}
\begin{minipage}[b]{\linewidth}\raggedright
\textbf{Período}
\end{minipage} & \begin{minipage}[b]{\linewidth}\raggedright
\textbf{Fase}
\end{minipage} & \begin{minipage}[b]{\linewidth}\raggedright
\textbf{Énfasis del gobierno}
\end{minipage} & \begin{minipage}[b]{\linewidth}\raggedright
\textbf{Incorporaciones}
\end{minipage} \\
\midrule\noalign{}
\endfirsthead
\toprule\noalign{}
\begin{minipage}[b]{\linewidth}\raggedright
\textbf{Período}
\end{minipage} & \begin{minipage}[b]{\linewidth}\raggedright
\textbf{Fase}
\end{minipage} & \begin{minipage}[b]{\linewidth}\raggedright
\textbf{Énfasis del gobierno}
\end{minipage} & \begin{minipage}[b]{\linewidth}\raggedright
\textbf{Incorporaciones}
\end{minipage} \\
\midrule\noalign{}
\endhead
\bottomrule\noalign{}
\multicolumn{4}{@{}l}{\rule{0pt}{3.2ex}\textbf{Fuente:} Elaboración propia.} \\
\endlastfoot
\textbf{Meses 1 a 12} & Levantamiento, arquitectura, desarrollo e integración de la Etapa 1; preparación de infraestructura & Líneas base, arquitectura base, ambientes, dependencias y adquisiciones. & Plan para la dirección en H1; arquitectura base en H2 (mes 4); Líder de Operación desde el mes 6; Líder de Implantación desde el mes 8. \\ \addlinespace
\textbf{Meses 13 a 15} & Marcha blanca de la Etapa 1 y desarrollo de la Etapa 2 & Preparación de soporte, observabilidad, continuidad y equipos; capacidad separada entre construcción y estabilización. & Comité de Operación desde el mes 13; sala de control diaria. \\ \addlinespace
\textbf{Meses 16 a 18} & Producción de la Etapa 1 y construcción de la Etapa 2 & Servicio de la Etapa 1 en producción y avance de la Etapa 2 sin reducir controles. & Cierre de la Etapa 1 en el mes 16. \\ \addlinespace
\textbf{Meses 19 y 20} & Etapa 1 en producción y marcha blanca de la Etapa 2 & Convivencia de versiones, autoridad única de datos compartidos, conciliación y retorno. & Sala de control de la Etapa 2. \\ \addlinespace
\textbf{Mes 21} & Producción de la Etapa 2 e inicio de la operación & Cierre de implementación y transferencia de ambos alcances a operación. & Informe de cierre de implementación. \\ \addlinespace
\textbf{Meses 21 a 56} & Operación & Niveles de servicio, incidentes, problemas, capacidad, continuidad, cambios evolutivos y transferencia de conocimiento. & El Comité de Operación asume el seguimiento ordinario; cierre contractual en el mes 56. \\
\end{longtable}

Antes de la producción de cada etapa, el Jefe de Proyecto reúne las evaluaciones de calidad, seguridad, datos, integración, implantación y operación; el Comité de Proyecto revisa la preparación y la contraparte técnica formaliza la conformidad. Si existe una condición bloqueante, se aplica el plan de recuperación o reversión sin sustituir el procedimiento de aceptación.

\subsection{T-9.8 Plantillas de gestión}\label{t9-8}

Las plantillas siguientes estandarizan los registros principales de la gestión. La Tabla~\mbox{\ref{tab:t9-pl-cambio}} corresponde a la solicitud de cambio.

\begin{longtable}[]{@{}
  >{\raggedright\arraybackslash}p{(\columnwidth - 4\tabcolsep) * \real{0.3000}}
  >{\raggedright\arraybackslash}p{(\columnwidth - 4\tabcolsep) * \real{0.7000}}@{}}
\caption{Plantilla de solicitud de cambio.}\label{tab:t9-pl-cambio}\\
\toprule\noalign{}
\begin{minipage}[b]{\linewidth}\raggedright
\textbf{Campo}
\end{minipage} & \begin{minipage}[b]{\linewidth}\raggedright
\textbf{Contenido}
\end{minipage} \\
\midrule\noalign{}
\endfirsthead
\toprule\noalign{}
\begin{minipage}[b]{\linewidth}\raggedright
\textbf{Campo}
\end{minipage} & \begin{minipage}[b]{\linewidth}\raggedright
\textbf{Contenido}
\end{minipage} \\
\midrule\noalign{}
\endhead
\bottomrule\noalign{}
\multicolumn{2}{@{}l}{\rule{0pt}{3.2ex}\textbf{Fuente:} Elaboración propia.} \\
\endlastfoot
\textbf{Identificador y fecha} & SC-NNN, fecha de registro y solicitante. \\ \addlinespace
\textbf{Origen} & Necesidad funcional, hallazgo técnico, incidencia u observación del cliente, con su vínculo. \\ \addlinespace
\textbf{Motivo y resultado esperado} & Qué se necesita cambiar y por qué. \\ \addlinespace
\textbf{Entregables y líneas base afectados} & Elementos de la EDT, requisitos, hitos y líneas base. \\ \addlinespace
\textbf{Análisis de impacto} & Alcance, plazo, costo, riesgo, niveles de servicio, arquitectura, interfaces, datos, seguridad, calidad, recursos, capacitación y operación. \\ \addlinespace
\textbf{Evaluaciones de frontera} & Líder de Datos, Encargado de Seguridad y contraparte de cumplimiento, cuando involucra datos u obligaciones del emisor. \\ \addlinespace
\textbf{Alternativas} & Opciones evaluadas y consecuencias de aprobar o rechazar. \\ \addlinespace
\textbf{Recomendación y decisión} & Recomendación del Comité de Proyecto, decisión del Comité Ejecutivo y acuerdo escrito si es contractual. \\ \addlinespace
\textbf{Efecto acumulado} & Efecto acumulado sobre el límite contractual, sin importes en la Oferta Técnica. \\ \addlinespace
\textbf{Ejecución y cierre} & Responsable, evidencia del resultado y fecha de cierre. \\
\end{longtable}

La Tabla~\mbox{\ref{tab:t9-pl-acta}} define el formato del acta de los cuatro comités.

\begin{longtable}[]{@{}
  >{\raggedright\arraybackslash}p{(\columnwidth - 4\tabcolsep) * \real{0.3000}}
  >{\raggedright\arraybackslash}p{(\columnwidth - 4\tabcolsep) * \real{0.7000}}@{}}
\caption{Plantilla de acta de comité.}\label{tab:t9-pl-acta}\\
\toprule\noalign{}
\begin{minipage}[b]{\linewidth}\raggedright
\textbf{Campo}
\end{minipage} & \begin{minipage}[b]{\linewidth}\raggedright
\textbf{Contenido}
\end{minipage} \\
\midrule\noalign{}
\endfirsthead
\toprule\noalign{}
\begin{minipage}[b]{\linewidth}\raggedright
\textbf{Campo}
\end{minipage} & \begin{minipage}[b]{\linewidth}\raggedright
\textbf{Contenido}
\end{minipage} \\
\midrule\noalign{}
\endhead
\bottomrule\noalign{}
\multicolumn{2}{@{}l}{\rule{0pt}{3.2ex}\textbf{Fuente:} Elaboración propia.} \\
\endlastfoot
\textbf{Instancia, sesión y fecha} & Comité, número de sesión, fecha y lugar o medio. \\ \addlinespace
\textbf{Asistentes} & Representantes autorizados de ambas partes y responsables de las materias tratadas. \\ \addlinespace
\textbf{Agenda y antecedentes} & Puntos tratados y documentos de entrada. \\ \addlinespace
\textbf{Decisiones} & Decisiones adoptadas, con su fundamento. \\ \addlinespace
\textbf{Acuerdos} & Acuerdo, responsable y plazo, vinculado a Jira. \\ \addlinespace
\textbf{Asuntos abiertos y escalados} & Asuntos sin resolver, motivo e instancia a la que se elevan. \\ \addlinespace
\textbf{Emisión} & Fecha de emisión, dentro de los dos días hábiles siguientes, y publicación en SharePoint. \\
\end{longtable}

La Tabla~\mbox{\ref{tab:t9-pl-leccion}} corresponde al registro de lecciones aprendidas, que alimentan el trabajo siguiente de cada frente.

\begin{longtable}[]{@{}
  >{\raggedright\arraybackslash}p{(\columnwidth - 4\tabcolsep) * \real{0.3000}}
  >{\raggedright\arraybackslash}p{(\columnwidth - 4\tabcolsep) * \real{0.7000}}@{}}
\caption{Plantilla de lección aprendida.}\label{tab:t9-pl-leccion}\\
\toprule\noalign{}
\begin{minipage}[b]{\linewidth}\raggedright
\textbf{Campo}
\end{minipage} & \begin{minipage}[b]{\linewidth}\raggedright
\textbf{Contenido}
\end{minipage} \\
\midrule\noalign{}
\endfirsthead
\toprule\noalign{}
\begin{minipage}[b]{\linewidth}\raggedright
\textbf{Campo}
\end{minipage} & \begin{minipage}[b]{\linewidth}\raggedright
\textbf{Contenido}
\end{minipage} \\
\midrule\noalign{}
\endhead
\bottomrule\noalign{}
\multicolumn{2}{@{}l}{\rule{0pt}{3.2ex}\textbf{Fuente:} Elaboración propia.} \\
\endlastfoot
\textbf{Identificador y frente} & LA-NNN, frente, etapa y fecha. \\ \addlinespace
\textbf{Contexto} & Actividad en que se observó: sprint, intervalo, ensayo de migración, ola, incidencia o comunicación. \\ \addlinespace
\textbf{Problema o práctica observada} & Qué ocurrió y qué efecto tuvo. \\ \addlinespace
\textbf{Respuesta aplicada} & Acción tomada en el momento. \\ \addlinespace
\textbf{Resultado} & Efecto de la respuesta. \\ \addlinespace
\textbf{Uso propuesto} & Ajuste de procedimiento, verificación adicional o preparación que se aplicará en el trabajo siguiente, con responsable. \\ \addlinespace
\textbf{Acceso} & Clasificación de la información, para no exponer datos o expedientes del emisor a personas no autorizadas. \\
\end{longtable}

\subsection{Referencias}\label{t9-referencias}

Las siguientes fuentes respaldan los procesos y controles citados en este formulario.

\noindent\hangindent=1.5em\hangafter=1 International Organization for Standardization. (2018). \emph{ISO 31000:2018. Risk management -- Guidelines}. ISO.

\noindent\hangindent=1.5em\hangafter=1 Kanban Guides. (2025, mayo). \emph{The Kanban Guide} (v2025.5). \url{https://kanbanguides.org/the-kanban-guide/2025.5/}

\noindent\hangindent=1.5em\hangafter=1 Ley N.º 21.719. (2024). \emph{Regula la protección y el tratamiento de los datos personales y crea la Agencia de Protección de Datos Personales}. Biblioteca del Congreso Nacional de Chile. \url{https://www.bcn.cl/leychile/Navegar?idNorma=1209272}

\noindent\hangindent=1.5em\hangafter=1 Prasad, V., Rajkumar, P., y Rastogi, S. C. (2006). \emph{Managing firmed fixed price projects using EVM: A case-study}. Project Management Institute. \url{https://www.pmi.org/learning/library/managing-fixed-price-projects-evm-8142}

\noindent\hangindent=1.5em\hangafter=1 Project Management Institute [PMI]. (2017). \emph{A guide to the project management body of knowledge (PMBOK® guide)} (6.ª ed.). Project Management Institute.

\noindent\hangindent=1.5em\hangafter=1 Prosci. (s.f.). \emph{The Prosci ADKAR model}. \url{https://www.prosci.com/methodology/adkar}

\ossFinalPage

\end{document}
